\documentclass[12pt,a4paper,fleqn]{article}

% Language setting
% Replace `english' with e.g. `spanish' to change the document language
\usepackage[english]{babel}

% automatically inserts a barrier after every \section (or \subsection, depending on the option), so that floats (figures/tables) cannot move past it.
\usepackage[section]{placeins}

% Set page size and margins
% Replace `letterpaper' with `a4paper' for UK/EU standard size
\usepackage[letterpaper,top=2cm,bottom=2cm,left=3cm,right=3cm,marginparwidth=1.75cm]{geometry}

% Useful packages
\usepackage{amsmath}
%\usepackage[doublespacing]{setspace}
\setlength{\parindent}{0pt}
\usepackage[flushmargin, hang]{footmisc}
%\usepackage[x11names]{xcolor}
\usepackage[bookmarks=false, colorlinks=true]{hyperref}
\AtBeginDocument{\hypersetup{citecolor=black}}
\AtBeginDocument{\hypersetup{linkcolor=black}}
\AtBeginDocument{\hypersetup{urlcolor=black}}
\usepackage{natbib}
\usepackage{bm}
\usepackage{amsmath}



\usepackage{amsfonts}
\usepackage{url}
\usepackage{float}
\usepackage{rotating}
\usepackage{lscape}
\usepackage[aboveskip=0pt,belowskip=0pt,labelfont=bf,textfont=bf,justification=justified,singlelinecheck=false]{caption} % font=small
\def\captionab#1#2{\caption[#1]{#1. \normalfont\footnotesize{#2}}}
\usepackage[section]{placeins} % For correct positioning of figures
\usepackage{multirow}
%\usepackage{todonotes}
\usepackage{authblk}
\usepackage{parskip}
\usepackage{booktabs}      
\usepackage{array}

\usepackage{graphicx}
\usepackage[justification=centering]{caption}
\usepackage{subcaption}
\usepackage{enumitem}

\usepackage{xcolor}
\usepackage{setspace}
\setstretch{1}
\usepackage[flushleft]{threeparttable}
\usepackage{bigstrut}

\usepackage{longtable} % Load longtable package
\usepackage{supertabular}

\usepackage{ragged2e}
\usepackage{lipsum}
%\usepackage{adjustbox}

\usepackage[ruled,vlined]{algorithm2e}
\usepackage{afterpage}

%\usepackage{arydshln} % Used for to draw dash-lines in array/tabular environments

% Make long url break at every letter in the alphabet
\expandafter\def\expandafter\UrlBreaks\expandafter{\UrlBreaks%  save the current one
      \do\a\do\b\do\c\do\d\do\e\do\f\do\g\do\h\do\i\do\j%
      \do\k\do\l\do\m\do\n\do\o\do\p\do\q\do\r\do\s\do\t%
      \do\u\do\v\do\w\do\x\do\y\do\z\do\A\do\B\do\C\do\D%
      \do\E\do\F\do\G\do\H\do\I\do\J\do\K\do\L\do\M\do\N%
      \do\O\do\P\do\Q\do\R\do\S\do\T\do\U\do\V\do\W\do\X%
      \do\Y\do\Z}


%\newcommand*\samethanks[1][\value{footnote}]{\footnotemark[#1]}

% Adapt \thanks symbol
\makeatletter
\def\@fnsymbol#1{\ensuremath{\ifcase#1\or *\or \mathsection\or \mathparagraph\or \|\or **\or \dagger\dagger \or \ddagger\ddagger \else\@ctrerr\fi}}
\makeatother

%%%%%%%%%%%%



\def\spacingset#1{\renewcommand{\baselinestretch}{#1}\small\normalsize}
\newgeometry{left=2.54cm, right=2.54cm, top=1cm, bottom=2.54cm}
\spacingset{1.2}	

\newcommand{\RNum}[1]{\uppercase\expandafter{\romannumeral #1\relax}}
\title{News-Based Multidimensional Political Instability in West and Central Africa\thanks{The authors thank the participants of the World Bank Seminars Series, Jan-Egbert Sturm, Alexander Rathke, and Samad Sarferaz for helpful discussions and comments.}} 

%Other option for the title:
%Real-Time Multidimensional Political Instability in West and Central Africa


\vspace{-0.25cm}
\author[$\ddagger$]{Philipp Kronenberg}
%\author[$\dagger$]{Kabinet Kaba}
\author[$\dagger$ $\lozenge$]{Pierre Mandon\thanks{\protect\url{pmandon@worldbank.org} [corresponding author].}}
\affil[$\ddagger$]{\small KOF Swiss Economic Institute, ETH Zurich, Switzerland}
\affil[$\dagger$]{\small World Bank Group, Washington, D.C., USA}
\affil[$\lozenge$]{\small CERDI \& University of Clermont Auvergne, Clermont-Ferrand, France}
\date{\today}

\begin{document}

\makeatletter
\let\origsection\section
\renewcommand\section{\FloatBarrier\origsection}
\let\origsubsection\subsection
\renewcommand\subsection{\FloatBarrier\origsubsection}
\makeatother

\maketitle

\begin{abstract}

\thispagestyle{empty}


\noindent Political instability remains a central obstacle to economic development in West and Central Africa. Despite extensive research, existing empirical work typically treats political instability as a one-dimensional construct and relies often on low-frequency, event-based datasets that are ill-suited for capturing short-term dynamics or identifying the timing of political shocks. This results in a limited understanding of how different forms of instability evolve and affect economic outcomes. This paper addresses these gaps by developing high-frequency, news-based indicators that separately track four dimensions of political instability. The results show that the four dimensions contain distinct information, spike at different historical moments, and display heterogeneous persistence and cross-border spillovers. In a case study for the Democratic Republic of the Congo, embedding these indicators in a VAR model %and linking them to firm-level data
%, we document different macroeconomic % and microeconomic
%effects of the four instability shocks on exchange rates, inflation, interest rates, and economic activity%, and sizable declines in firm sales and employment
%.\\
reveals heterogeneous transmission on the economy: violence shocks resemble localized risk and supply disruptions; protest shocks generate short-term risk increases; within-regime instability maps to slower, more persistent uncertainty effects; and regime instability behaves like a broad confidence shock with demand restraint.\\

\noindent \textbf{JEL Classification:} C43, C82, E32, F51, F52, N47 \\
\noindent \textbf{Keywords}: Political instability, developing countries, VAR, news sentiment, Africa


\end{abstract}
\clearpage

\newpage

\pagenumbering{arabic}
\setcounter{page}{1}

\newgeometry{left=2.54cm, right=2.54cm, top=2.54cm, bottom=2.54cm}
\spacingset{1.84}



\section{Introduction}

Political instability remains a central obstacle to economic development in West and Central Africa. The region has experienced more coups d'état than any other since independence, with a pronounced resurgence since 2020: military takeovers in Mali (2020, 2021), Guinea (2021), Burkina Faso (2022, twice), Niger (2023), and Gabon (2023) have reversed democratic gains and disrupted regional integration \citep{mcgowan2003}. Frequent episodes of coups, mass protests, and contested regime turnover have shaped governance structures, undermined policy continuity, and weighed on long-term economic performance. Political instability erodes investor confidence, weakens state institutions, and increases the risk of violent conflict, thereby constraining growth and poverty reduction \citep{Alesina1996, VERDUGO2013, aisen2013, collier2002, CHEN1996, DARBY2004, DEVEREUX1998}. Given the prevalence of fragile political systems and cross-border spillovers -- especially in the Sahel and the Congo Basin -- systematic measurement of political instability is essential for early-warning systems and evidence-based policy design \citep{goldstone2010, hegre2017}.

Despite extensive research on political instability and its economic consequences, the empirical literature faces two persistent challenges. The first is conceptual: much of the empirical work treats political instability as a one-dimensional phenomenon, typically proxied by composite indices or a narrow subset of events such as coups, assassinations, or riots \citep{Feierabend1966, Rummel1966, Tanter1966, Alesina1996, Perotti1996}. This aggregation masks the multidimensionality of political instability highlighted by \cite{JONGAPIN2009}, who identifies four distinct dimensions: politically motivated violence, mass civil protest, instability within the political regime, and instability of the political regime. These dimensions need not move together and may have different economic implications. For example, violent conflict is likely to depress investment and destroy productive capacity, whereas intra-regime turnover or heightened parliamentary contestation may, in some contexts, strengthen accountability. When diverse forms of political instability are collapsed into a single index, it remains unclear which specific aspects drive macroeconomic % and microeconomic 
outcomes.

The second challenge is measurement. The dominant approach relies on low-frequency, event-based, or survey-based indicators. Because the most widely used political instability and conflict datasets record the onset of discrete events at annual or multi-year frequency, they are valuable for long-run cross-country comparisons but are less appropriate for capturing short-term fluctuations, forecasting, or tracing the transmission of political shocks in macroeconomic models. Their categorical nature and retrospective coding make it difficult to identify the precise timing and intensity of political shocks. As a result, empirical research linking political instability to economic outcomes -- such as investment, inflation, or growth -- often suffers from weak identification.

To address these limitations, recent work has begun to exploit high-frequency textual sources to construct continuous, news-based sentiment indicators that quantify otherwise latent risks. \cite{baker2016} develop a widely used index of economic policy uncertainty; \cite{AHIR2022} extend this approach to construct a world uncertainty index; \cite{caldara2022} measure geopolitical risk; and \cite{BARRETT2022} track social unrest. In contrast to traditional event datasets, news-based indicators are continuous, timely, and easily scalable across countries, making them particularly well suited for VAR models, state-space models, and forecasting frameworks that require regular and up-to-date inputs \citep{shapiro2022, ahir2023financial, APRIGLIANO2023791, fraiberger2021media, HUANG2019}. However, none of the existing news-based measures capture the multidimensional character of political instability.

In this paper, we address these gaps by developing news-based political instability sentiment indicators (NBS) for four dimensions of instability, following the classification of \cite{JONGAPIN2009}. Using international newspaper archives from Dow Jones Factiva, we construct monthly indices that separately track politically motivated violence, mass civil protest, instability within the political regime, and instability of the political regime for 17 West and Central African countries over the period 1990--2024. Our methodology builds on the text-search and validation strategies of \cite{caldara2022} and \cite{BARRETT2022}, adapting them to a multidimensional concept of political instability and to the institutional context of low- and lower-middle-income countries. Whereas \cite{DIRKS2024102586} examine political instability and growth in 34 advanced economies using the WGI Political Stability index, we extend this analysis to developing countries where instability is more prevalent and its economic consequences potentially more severe.

Our results show that the four dimensions of political instability contain distinct signals. The pairwise correlations between dimensions indicate that they capture related but separable phenomena. The indices spike at different historical events: the political violence index peaks during the Mali Tuareg rebellions (1990--1996, 2012) and the Boko Haram insurgency (2009--present); the mass protest index spikes during the Senegal Y'en a Marre movement (2011--2012) and the Nigeria EndSARS protests (2020); and the regime instability index surges during the post-2020 coup wave. We document meaningful cross-border correlation, especially for political violence and mass civil protests, giving evidence that instability shocks might propagate to neighboring countries, while instability within the regime and instability of the regime shocks are rather country-specific. In a case study of the Democratic Republic of the Congo, embedding our indicators in a VAR model using high-frequency data, we find that our four political instability shocks affect macroeconomic variables in different ways: while political violence shocks resemble localized risk and supply disruptions, mass civil protests create short-term risk increases, within-regime instability maps slower, more persistent uncertainty, and regime instability behaves like a broad confidence shock with demand restraint. %Using firm-level panel data from the World Bank Enterprise Surveys, we also document sizable effects of instability shocks on sales and employment, highlighting how political risk transmits to the microeconomic margin.

This paper makes four contributions to the political instability literature. First, we develop high-frequency, news-based indicators of political instability tailored to West and Central African countries, where conventional data sources are often sparse, delayed, or subject to political interference. This approach complements the growing use of alternative data sources -- such as satellite imagery and nighttime lights \citep{henderson2012, DONALDSON2016} -- to measure economic activity in data-scarce environments. Second, we operationalize the multidimensional framework of \cite{JONGAPIN2009} in a news-based setting, allowing for a more nuanced decomposition of political shocks than single-index approaches. Third, by linking our indicators to %both
macroeconomic aggregates % and firm-level outcomes, 
we provide one of the first systematic assessments of how different types of political instability propagate through the economy in developing countries, complementing existing work focused on advanced economies \citep{DIRKS2024102586}. Fourth, our framework generates indicators that can be updated in near real-time, contributing to the emerging literature on early-warning models of conflict and instability \citep{goldstone2010, hegre2017}.

The remainder of the paper is structured as follows. Section~\ref{sec:construction} describes the conceptual framework, reviews existing measurement approaches, and explains how we construct our news-based indices. Section~\ref{sec:properties} documents the properties of the four dimensions, including their historical evolution, cross-country patterns, and validation against benchmark indicators. Section~\ref{sec:empirical_analysis} presents the empirical applications: VAR models linking instability shocks to macroeconomic outcomes% and firm-level panel regressions examining effects on sales and employment
. Section~\ref{sec:conclusion} concludes.


\section{Construction of the Political Instability Index}\label{sec:construction}

This section presents the conceptual framework underlying our measurement approach, reviews existing data sources, describes the news data and keyword search methodology, and defines the index construction.


\subsection{Conceptual Framework}\label{sec:concept}

Political instability refers to the uncertainty and unpredictability within a country's political system, often resulting from frequent government changes, weak institutions, civil unrest, or conflicts. Its consequences are far-reaching, hindering economic growth, discouraging investment, and, in extreme cases, leading to state failure or prolonged civil conflict. In empirical research, political instability has traditionally been treated as a one-dimensional concept \citep{Alesina1996}. However, measuring political instability in a multidimensional way can yield new insights \citep{Perotti1996, Rummel1966, Tanter1966, Feierabend1966, Morrison1971, hibbs1973mass}.

We follow \cite{JONGAPIN2009}, who uses principal component analysis on 25 political instability indicators to identify four distinct dimensions: politically motivated violence, mass civil protest, instability within the political regime, and instability of the political regime. Each dimension captures a different facet of political disruption. Politically motivated violence encompasses armed conflict, insurgency, terrorism, and state repression -- the physical contestation of power through force. Mass civil protest refers to collective action by citizens, including demonstrations, strikes, and riots, that challenge authority without necessarily involving armed combat. Instability within the political regime encompasses elite fragmentation, legislative gridlock, coalition breakdown, and contested succession -- the internal dysfunction of governing institutions. Instability of the political regime refers to fundamental challenges to regime continuity, including coups d'état, constitutional crises, and forced leadership changes.

This multidimensional framework is especially relevant for West and Central Africa, where all four dimensions are empirically salient. The Sahel has experienced sustained jihadist insurgencies since 2012, with violence spreading across Mali, Niger, Burkina Faso, and Chad. Simultaneously, mass protest movements have emerged along the coast: Senegal's Y'en a Marre (2011--2012), Nigeria's EndSARS (2020), and Togo's anti-constitutional protests (2017--2019). Elite instability characterizes long-ruling regimes in Chad, Cameroon, and Equatorial Guinea, where succession uncertainties generate persistent political tension. And the post-2020 coup wave -- Mali (2020, 2021), Guinea (2021), Burkina Faso (2022, twice), Niger (2023), Gabon (2023) -- demonstrates that regime instability remains a defining feature of the region. Treating these diverse phenomena as a single index would obscure their distinct drivers and economic implications.


\subsection{Existing Approaches to Measurement}\label{sec:existing}

Existing data sources for measuring political instability can be grouped by their data generation process: event-based datasets, expert- and perception-based indicators, interview-based surveys, and media- and text-based approaches. Each offers unique analytical strengths but also distinct limitations.

Event-based datasets record observable political actions, typically coded from official records, historical archives, or news reports. The Database of Political Institutions \citep{Beck2001} provide annual, country-level data on institutional characteristics since 1975, the Polity V Project \citep{marshall2020} evaluates regime types and transitions since 1800, and the Cross-National Time-Series archive (CNTS) \citep{banks2013cross} records conflict events as well as political indicators. The Powell and Thyne coup dataset \citep{POWELL2011} records successful and failed coups since 1950. The Uppsala Conflict Data Program (UCDP) \citep{sundberg2012} documents armed conflicts, organized violence, and peace processes since 1946. The Armed Conflict Location \& Event Data Project (ACLED) \citep{raleigh2010} provides high-frequency, geo-referenced data on political violence and protest in Africa since 1997. The World Handbook of Political Indicators IV (WHIV) \citep{jenkins2012} records protests, strikes, and coups alongside cooperative actions between 1979--1992. These datasets offer structured, cross-nationally comparable evidence over long time horizons, but they capture only realized events, not latent tensions, and their low temporal frequency limits their usefulness for short-run macroeconomic analysis.

Expert- and perception-based indicators convert qualitative assessments into quantitative measures of stability. The International Country Risk Guide (ICRG) \citep{howell2011} produces monthly political risk scores based on qualitative analyst judgment since 1984. The Worldwide Governance Indicators (WGI) \citep{kaufmann2011} aggregate over 30 sources into six governance dimensions for over 200 countries since 1996. These indicators enable global comparison but are subjective, slow-moving, and may lag actual political developments.

Interview-based surveys capture citizens' perceptions and attitudes. The Gallup World Poll \citep{tortora2010} is a globally standardized survey conducted
annually since 2005, covering well-being, trust in institutions, economic perceptions, and social attitudes across more than 140 countries. The Afrobarometer \citep{mattes2008} collects nationally representative surveys on democracy, governance, and social trust across over 40 African countries in multi-year rounds. These data provide insight into latent instability, such as distrust or dissatisfaction, before it becomes manifest, but are constrained by temporal gaps and limited coverage in fragile states. 

Media- and text-based indicators represent a more recent approach. The Global Database of Events, Language, and Tone (GDELT) \citep{ward2013} applies machine learning to extract events from global news updated every 15 minutes. The Reported Social Unrest Index \citep{BARRETT2022} tracks unrest-related keywords in 18 English-language news outlets for 131 countries since 1985. It quantifies media attention to protests and riots, offering a synthetic indicator of societal
tension. The World Uncertainty Index \citep{AHIR2022} measures economic and political uncertainty based on the frequency of ``uncertainty'' terms in Economist Intelligence Unit reports since 1952. These indicators offer high-frequency, global coverage but capture media salience rather than objective instability, and they typically aggregate instability into a single dimension.

Our approach builds on this media-based tradition while addressing key limitations. We construct separate indices for each of the four \cite{JONGAPIN2009} dimensions rather than collapsing instability into a single measure. We apply transparent, replicable keyword searches tailored to the West and Central African context. And we validate our indices against established benchmarks to assess convergence validity.


\subsection{Data}\label{sec:data}

We use newspaper data from Dow Jones' Factiva news aggregator, which provides more than 32,000 sources in 28 languages, and 600 newswires. We focus on countries belonging to the Central African Economic and Monetary Community (French: Communauté économique et monétaire de l'Afrique centrale, or CEMAC) (Cameroon, Central African Republic, Chad, Equatorial Guinea, Gabon, Republic of Congo) and West African Economic and Monetary Union (WAEMU) (Benin, Burkina Faso, Côte d'Ivoire, Guinea-Bissau, Mali, Niger, Senegal, Togo) monetary unions, which provide a high degree of institutional and macroeconomic comparability through shared monetary policy, fiscal frameworks, and administrative traditions. To capture broader security dynamics relevant to political instability, we extend the analysis to the French-speaking Sahel (Mali, Niger, Burkina Faso, Chad, Senegal, Mauritania), where countries face similar exposure to political violence, military interventions, and cross-border insurgencies, as well as the non-French-speaking Ghana and Nigeria, which are the region’s largest economic and political actors. Countries outside these groupings -- such as Gambia, Guinea, Liberia, Sierra Leone, Rwanda, and Burundi -- are excluded due to differing colonial legacies, monetary regimes, institutional structures, or conflict dynamics, which would reduce comparability and complicate interpretation. Figure \ref{fig:africa_map} in the Appendix shows the relevant countries colored in red on a map. 


% Right now map is place holder. For later: To show the relevance of the topic: Could adjust the map to visualize any type of political instability to show the importance of the measure in the region... 

We selected the most important global news sources for African countries according to the Global Media Index for Africa 2024 which was created in collaboration with The Africa Center, Africa No Filter, and the University of Cape Town. Table \ref{tab:sources_summary} provides a summary of the used news sources. The news sources are all in English and are updated either daily or weekly.

Our motivation for the choice of news sources is to create a sample that has high coverage and quality, and minimal bias for the countries under study. The coverage and quality of local news sources may vary strongly. Local newspapers may have different use of language or cover only very specific topics. Further, they may cover rather topics of local interest, whereas we are rather interested in important events and news that might have an impact on the economy as a whole, therefore, focusing on the perception of events from an international perspective should filter out less relevant news. Finally, political biases cannot be ruled out for local news sources. This is particularly relevant for political instability shocks, which are mostly driven by political actions of the leading political party or other influential political groups. For this reason, to provide a consistent and reliable measure of news reporting, we restrict our sample to major English-language newspapers and networks, including news sources from all over the world. To obtain a balanced sample, we included not only American newspapers, which have the most extensive worldwide coverage, but also European newspapers, due to their historical importance, as well as non-Western newspapers, due to their growing influence in sub-Saharan Africa. 


\begin{table}[h]
   \caption{Summary of Newspaper Sources}
   \label{tab:sources_summary}
   \centering
   \scriptsize
   \renewcommand{\arraystretch}{1.5} % <-- increases vertical spacing
   \resizebox{\linewidth}{!}{%
      \begin{tabular}{llllll}
         \hline \addlinespace[0.3em]
         \hline
         Source & Language & First Publication & Frequency & Availability & Type \\ 
         \addlinespace[0.3em]
         \hline
         Agence France Presse & English & 1991-09-09 & Daily & Continuously updated & Newswires \\ 
         Associated Press Newswires & English & 1985-12-03 & Daily & Continuously updated & Newswires \\ 
         BBC Monitoring Newsfile & English & 1998-09-25 & Daily & Continuously updated & Newswires \\ 
         DW English & English & 2006-06-30 & Daily & Continuously updated & Media Transcripts \\ 
         The Economist & English & 1981-01-04 & Weekly & 1 day before publication & Various \\ 
         The Guardian & English & 1981-01-02 & Daily & Continuously updated & Newspapers \\ 
         Al Jazeera English & English & 2007-01-05 & Daily & Continuously updated & Newswires \\ 
         Reuters News & English & 1987-05-27 & Daily & Continuously updated & Newswires \\ 
         Voice of America & English & 2002-01-01 & Daily & Same day as publication & Various \\ 
         The Wall Street Journal & English & 1979-06-13 & Mo-Sa & Start of local business day & Newspapers \\ 
         The Washington Post & English & 1977-01-01 & Daily & Start of local business day & Newspapers \\ 
         Xinhua News Agency & English & 1989-03-01 & Daily & Continuously updated & Newswires \\ 
         \hline
         \hline
      \end{tabular}
   }%
   \vspace{0.5em}
   \begin{minipage}{\textwidth}
   {\footnotesize \flushleft \textit{Notes:} This table summarizes the various news sources, their availability, frequency, and types.}
   \end{minipage}
\end{table}


To have a large enough sample but at the same time enough articles, we decided to start our analysis in 1990, as most of the news sources had already started publishing by then. Only Deutsche Welle (DW English), Al Jazeera English, and Voice of America start publishing later. However, we still include them to have a broader and more international coverage.

To measure a consistent signal over time, we need to control for trends in the overall articles of our news sources. In Figure \ref{fig:z_plot_all_cnt} in the Appendix, we illustrate the total number of articles per month for the countries relevant to our analysis. The figure shows that the number of articles has an upward trend from the 1990s until 2000. After that, the overall number of articles fluctuates but stays rather stable without a general trend. A look at the number of articles by country in Figure \ref{fig:y_plot_all_cnt} shows that there are some differences between the countries. Larger countries generally have more published articles, and certain crisis periods are more strongly pronounced than in the overall figure. The trend in the overall number of articles is primarily driven by Nigeria and the Democratic Republic of Congo, as they have much larger news coverage. For this reason, it makes sense to account for country-specific trends rather than an overall trend.


\begin{figure}[h]
\centering
\caption{Number of Articles per Month}
\includegraphics[width=\textwidth]{Figures/country_normalization_reference_counts.pdf}
\label{fig:y_plot_all_cnt}
\begin{minipage}{\textwidth}
 % Adds a bit of vertical space between the figure and the notes
    {\footnotesize{\textit{Notes:} The figure shows the number of monthly articles excluding all non-articles, lifestyle content, and routine news (see a full list in Table \ref{tab:search_terms}) for the countries included in the study (Benin, Burkina Faso, Cameroon, Central African Republic, Chad, Côte d'Ivoire, Democratic Republic of the Congo, Equatorial Guinea, Gabon, Ghana, Guinea-Bissau, Mali, Mauritania, Niger, Nigeria, Senegal, Togo) in black as well as the twelve-month rolling-mean in red.}}
\end{minipage}
\end{figure}

\subsection{Keyword Search}\label{sec:keywords}

To create an index of political instability, we perform a keyword search on the news sources. A large literature related to finance has shown that attention measures, such as Google searches or newspaper attention, have predictive power for stock market volatility predictions and that attention measures often have equal or higher predictive power than sentiment measures (see \cite{AUDRINO2020} and references therein). Several studies in macroeconomic forecasting have also shown the predictive power of such attention measures for macroeconomic variables such as policy-uncertainty by \cite{baker2016}, consumer sentiment by \cite{doms2004}, and unemployment and monetary policy by \cite{fisher2022}. We build on this literature and argue that the simple attention to the topic is enough to extract a signal. In the case of political instability, the sentiment is mostly negative and thus does not provide relevant additional information to the attention measure. 


For the article search, we use inclusive and exclusive requirements. Table \ref{tab:search_terms} in the Appendix shows an overview of the applied search terms and restrictions. Inclusive requirements define certain text strings or tags that need to be included in the article so that it is related to a specific event. In this way, we reduce false negatives (type II error). Restrictions on location and subject using tags define the country related to the article and the general topic of the article. Exclusive requirements aim to avoid false positives (type I error). Excluding certain strings helps avoid certain words that may be misleading or have different meanings in different contexts. For this purpose, we follow the string lists of \cite{caldara2022} and \cite{BARRETT2022}. Furthermore, subject tags are used to exclude non-relevant topics from the analysis. We exclude non-articles, lifestyle content, and routine news (weather, sports, images, etc.). 

Following \cite{JONGAPIN2009}, we try to extract the different dimensions of political instability as they are distinct from each other, and relying only on one specific dimension would be too restrictive. The dimensions are ``Politically motivated Violence'', ``Mass Civil Protests'', ``Instability within the Regime'', and ``Instability of the Regime''. For each of the dimensions, we define our own keyword searches with different inclusive requirements. The dimension of ``Politically motivated Violence'' is mostly correlated with indicators measuring guerrilla warfare, revolutions, civil wars, internal conflicts, ethnic tensions, assassinations, religious tensions, and other civil conflicts. As a result, keywords include terms related to political violence and warfare. The dimension of ``Mass Civil Protests'' is mostly correlated with the occurrence of riots, strikes, and demonstrations. The dimension of ``Instability within the Regime'' is mostly correlated with indicators describing the fractionalization and polarization of the political regime, as well as the years of the largest party in office, the number of elections, changes in the chief executive, and the number of veto players who leave office. As the analysis of \cite{JONGAPIN2009} has shown, those indicators are directly linked to the level of democracy or autocracy. As a result, we defined the keywords capturing how democracy is discussed in the media. Lastly, the dimension of ``Instability of the Regime'' is mostly correlated with changes in the chief executive, cabinet changes, major constitutional changes, coups d'état, the number of veto players who drop from office, political regime changes, major government crises, and measures of government stability. Thus, we tried to capture with the keywords any changes in the political regime.



The summary statistics reveal how distributional properties vary by country. Table \ref{tab:summary_statistics} in the Appendix shows the number of political instability articles for each month and country. It shows that the coverage varies significantly among countries. The distribution is right-skewed, with certain months having many more articles than the average. Those months indicate specific events. Often, those events are related to topics in which we are interested. While countries with stable, high coverage provide more information from which to extract a signal, the probability of false negatives (type II error) is low; however, the probability of false positives (type I error) is higher, as the data includes more irrelevant information from which we need to extract the signal relevant to us. On the other hand, for countries with small media coverage, it might be difficult to extract any signal, thus the risk of false negatives is higher; while the articles published in the international press are likely to be highly relevant, resulting in a low probability of false positives (type I error).  Table \ref{tab:summary_statistics_all} in the Appendix provides an overview of the different dimensions of political instability. As we conduct more search queries for political violence, which include more detailed keyword searches, the number of relevant articles is, by design, higher. However, countries that have a high (low) number of articles in one dimension also tend to have a high (low) number of articles in the other dimensions.



\subsection{Definition of the Index}\label{sec:index}

The Political Instability Index follows the definition of \cite{BARRETT2022}:
\begin{equation}
    IndexA_{it} = \frac{x_{it}}{\frac{1}{12} \sum_{j=1}^{12} z_{t-j}} \times \frac{100}{\bar{x}_i/\bar{z}},
\end{equation}
\vspace{-7em}

\begin{equation}
    IndexB_{it} = \frac{x_{it}}{\frac{1}{12} \sum_{j=1}^{12} y_{i,t-j}},
\end{equation}
where $x_{it}$ represents the number of articles about political instability in country $i$ at time $t$, while $z_t$ denotes the total number of contemporary articles at time $t$. Similarly, $y_{it}$ refers to the number of contemporary articles in country $i$ at time $t$. The term $\bar{x}_i$ corresponds to the average of $x_{it}$ over the whole period for country $i$, and $\bar{z}$ represents the average of $z_t$ over the whole period and country. As a result of the normalization, the two indices have different characteristics. 

Index A is defined as the fraction of relevant articles from country $i$ relative to the total number of articles for all countries included in the analysis. This enables us to normalize and correct for common trends in newspaper coverage in the region, while being more stable than examining the country-specific overall newspaper coverage; at the same time, it is not influenced by developments in newspaper coverage in other regions of the world, which may strongly deviate from those in the region relevant to the analysis. By scaling the index to a mean of 100 and correcting for country-specific media attention, the measure becomes robust over time but does not provide a measure of intensity across countries. 

The Index B is defined as the share of articles about political instability from the total number of contemporary articles for this country. By normalizing with a country-specific rolling mean, this index becomes less volatile and smooths observations over time, giving small events during periods of low newspaper coverage a higher weight, while strong signals during periods of high coverage are down-weighted. This introduces a varying change in the sensitivity of the index over time. Problems with Index B arise during times of very low coverage in the country, such that small events become very important, as well as the endogeneity of general news coverage to political instability, since public interest in the country generally increases during such periods. 



\section{Properties of the Index}\label{sec:properties}

This section documents the properties of our news-based political instability indices. We first examine each dimension separately, using historical narratives to validate that the indices spike around known instability events. We then analyze cross-dimensional patterns and report correlations with established benchmark indicators to assess convergent validity.


\subsection{Politically Motivated Violence}\label{sec:political_violence}

The indicator in Figure \ref{fig:country_count_norm_A_political_violence_NBS_A} primarily captures acts of collective or organized violence with explicit political motives -- including coups, rebellions, insurgencies, separatist conflicts, and armed repression. It reflects both violent challenges to the state and violent state responses to civilians. Across the region, the measure spikes during major insurgencies or military interventions. In Mali, large increases align with the Tuareg rebellions (1990-1996, 2006, 2012) and the jihadist insurgency from 2012 onward. Niger exhibits similar peaks during the Tuareg rebellions of the late 1990s and 2007, and again after 2010 as Boko Haram violence spread into the Diffa region. Nigeria shows repeated surges: the 1990 Gideon Orkar coup attempt, the Niger Delta militancy (1998-2009), Boko Haram’s insurgency since 2009, and periodic communal conflicts in Kaduna and Jos. Chad and the Central African Republic record near-continuous violent instability linked to civil wars, coups, and rebel incursions, particularly the 2003 and 2013 regime changes in Bangui. Côte d’Ivoire’s spikes correspond to the 2002-2007 civil war and the 2010-2011 post-election crisis. Smaller but distinct episodes of violence appear in Senegal’s Casamance insurgency, Guinea-Bissau’s coups (1998, 2003, 2012), and Togo’s 1993-94 and 2005 electoral crises. Equatorial Guinea, Gabon, and Mauritania record shorter bursts, mainly during attempted coups or transitions (e.g., Mauritania 2005, 2008). Overall, this measure traces the physical contestation of power, showing pronounced regional clustering in the Sahel and the Gulf of Guinea.\\

\begin{figure}[!t]
\centering
\caption{Politically motivated Violence Index per Country}
\includegraphics[width=\textwidth]{Figures/nbs_a_political_violence_by_country.pdf}
\label{fig:country_count_norm_A_political_violence_NBS_A}
\begin{minipage}{\textwidth}
 % Adds a bit of vertical space between the figure and the notes
    {\footnotesize{\textit{Notes:} This figure displays the A Index for Politically motivated Violence from 1990 to 2025.}}
\end{minipage}
\end{figure}

\subsection{Mass Civil Protest}\label{sec:mass_civil_protest}

The indicator in Figure \ref{fig:country_count_norm_A_mass_civil_protest_NBS_A} measures nonviolent or low-intensity collective action, demonstrations, strikes, and protests challenging authority, policy, or economic conditions. It reflects the mobilization capacity of civil society and opposition movements. Since the 1990s, protests have become a defining feature of political life across West and Central Africa, often accompanying liberalization and democratic backsliding. Senegal’s 2011 Y’en a Marre and 2021-2024 anti-third-term protests represent some of the largest and most sustained waves of civic mobilization in the region. Nigeria’s 2012 Occupy Nigeria and 2020 EndSARS protests marked milestones in digital youth activism, while the 2024 EndBadGovernance movement continued that trend. Togo and Benin saw repeated mobilizations over constitutional reform and election fairness (especially in 2017-2018 in Togo and 2019 in Benin). Niger and Mali both experienced protest peaks preceding or following coups, often over governance failures or security crises. In Ghana and Côte d’Ivoire, protest waves tended to accompany electoral controversies rather than systemic collapse, while Mauritania’s and Equatorial Guinea’s demonstrations remained localized and quickly contained. The indicator, therefore, captures societal expressions of discontent and regime contestation without outright violence, often preceding or accompanying elite instability.\\

\begin{figure}[!t]
\centering
\caption{Mass Civil Protest Index per Country}
\includegraphics[width=\textwidth]{Figures/nbs_a_mass_civil_protest_by_country.pdf}
\label{fig:country_count_norm_A_mass_civil_protest_NBS_A}
\begin{minipage}{\textwidth}
 % Adds a bit of vertical space between the figure and the notes
    {\footnotesize{\textit{Notes:} This figure displays the A Index for Mass Civil Protest from 1990 to 2025.}}
\end{minipage}
\end{figure}

\subsection{Instability of the Political Regime}\label{sec:instability_of_the_regime}

The indicator in Figure \ref{fig:country_count_norm_A_instability_of_regime_NBS_A} measures systemic breakdowns of governing authority -- episodes when the existing regime loses or transfers power outside constitutional norms, typically through coups, revolutions, or the collapse of control. It therefore marks thresholds of discontinuity in political order. The most prominent spikes occur in Mali (1991, 2012, 2020, 2021), Niger (1996, 1999, 2010, 2023), and Burkina Faso (2014 and 2022 coups), each producing new transitional regimes. Chad’s 2021 succession crisis, following President Idriss Déby’s death in combat, fits this category as an irregular power transfer. Mauritania saw clear regime breakdowns in 2005 and 2008, while Côte d’Ivoire’s civil war (2002–2011) and Guinea-Bissau’s repeated coups mark chronic systemic fragility. Togo’s hereditary handover in 2005 and Gabon’s 2023 military ouster of Ali Bongo illustrate the continuity of authoritarian regimes under new leadership. More stable democracies, such as Ghana, Benin, and Senegal, record minor, transient increases linked to electoral turnovers rather than collapses, underscoring regime durability despite protests. The Democratic Republic of Congo’s recurring armed conflicts (1996–2003; 2018 transition) represent protracted systemic crises. This indicator, in sum, traces episodes where formal political order itself is reconstituted, offering a long-term measure of regime resilience or fragility across the region.\\


\begin{figure}[!t]
\centering
\caption{Instability of the Regime Index per Country}
\includegraphics[width=\textwidth]{Figures/nbs_a_instability_of_regime_by_country.pdf}
\label{fig:country_count_norm_A_instability_of_regime_NBS_A}
\begin{minipage}{\textwidth}
 % Adds a bit of vertical space between the figure and the notes
    {\footnotesize{\textit{Notes:} This figure displays the A Index for Instability of the Regime from 1990 to 2025.}}
\end{minipage}
\end{figure}


\subsection{Instability Within the Political Regime}\label{sec:instability_within_the_regime}

The indicator in Figure \ref{fig:country_count_norm_A_instability_within_regime_NBS_A} reflects internal erosion of cohesion among ruling elites. While in democratic states the indicator is mainly driven by discussions about national elections, in less stable states it also captures elite-level fractures, power struggles, and realignments that occur without full regime overthrow -- including coups within ruling coalitions, purges, constitutional manipulation, and contested successions. Mali and Burkina Faso show recurrent within-regime instability through successive coups and reshufflings (2012, 2020, 2021 in Mali; 2014, 2022 in Burkina Faso). Mauritania exhibits repeated military interventions (2005, 2008) followed by controlled transitions, indicating elite self-correction mechanisms rather than full collapse. Chad’s pattern of rebellions and leadership succession (Déby’s death in 2021, followed by his son’s takeover) fits this category as well. Niger’s history alternates between elite fissures and restoration -- from the 1996 and 1999 coups to the 2023 Presidential Guard takeover. Nigeria’s numerous intra-military tensions during the 1990s, constitutional disputes under Obasanjo (2006-2007), and recent party infighting under Tinubu similarly contribute. Senegal, Togo, and Gabon experienced within-regime shifts via constitutional reforms aimed at perpetuating incumbents, producing detectable though milder spikes. Equatorial Guinea and Cameroon show persistent low-grade instability tied to opaque succession and internal elite management. Overall, this indicator detects intra-regime adaptation and competition, reflecting how regimes evolve under pressure while maintaining continuity.\\

\begin{figure}[!t]
\centering
\caption{Instability within the Regime Index per Country}
\includegraphics[width=\textwidth]{Figures/nbs_a_instability_within_regime_by_country.pdf}
\label{fig:country_count_norm_A_instability_within_regime_NBS_A}
\begin{minipage}{\textwidth}
 % Adds a bit of vertical space between the figure and the notes
    {\footnotesize{\textit{Notes:} This figure displays the A Index for Instability within the Regime from 1990 to 2025.}}
\end{minipage}
\end{figure}


\subsection{Cross-Dimensional and Cross-Country Patterns}\label{sec:cross_patterns}

To assess the validity of our news-based indices, we examine their correlation with established alternative measures of political (in-)stability. Table~\ref{tab:mean_correlation_selected_events_ordered} in the Appendix reports pairwise correlations between our four NBS dimensions and selected benchmark indicators.

The correlations between the NBS dimensions are positive but well below unity, supporting the view that the four dimensions capture separable phenomena with potentially distinct economic implications. Looking at the correlation with selected benchmark indicators, several patterns emerge. First, each NBS dimension correlates most strongly with its conceptually closest benchmark: the political violence index correlates with internal conflicts, civil war, coups, violence, ethnic violence, battles, and explosions, terrorism; the mass protest index correlates with protests, riots, corruption, and political sanctions; and the regime instability index correlates with government crises, and negatively with years in office, government effectiveness, the party in power, and the size of the major party. Second, there are benchmark measures that correlate with multiple dimensions as they measure political instability in a broader sense, such as civil unrest, internal conflicts, or violence. Third, all four dimensions correlate negatively with the WWGI political stability score, as expected. These patterns suggest that our news-based indices capture meaningful variations in political instability that align with established political (in-)stability measures while providing the additional advantages of higher frequency and continuous measurement.

Cross-country correlations reveal geographic clustering. The political violence index shows the highest spatial correlation, particularly among Sahel countries (Mali, Niger, Burkina Faso), where jihadist insurgencies have spread across borders, and among Great Lakes countries (DRC, CAR, Chad) with interconnected conflict dynamics. The mass civil protest index shows lower cross-country correlation, as protest waves tend to be driven by country-specific political calendars. Similarly, the within-regime instability index is primarily linked to political election processes and country-specific political conflicts. Regime instability shows episodic regional clustering, most notably during the post-2020 coup wave, when military takeovers in one country appeared to embolden similar actions in neighboring countries. 



\subsection{Evolution over Time}\label{sec:evolution_over_time}

To assess the overall evolution of political instability, we construct a composite index that captures deviations from the historical mean for each dimension. The composite index in Figure \ref{fig:country_count_by_topic_stacked_avg} is based on an equally weighted average of each country’s Index B across the four dimensions of political instability -- political violence, mass civil protest, instability within the regime, and instability of the regime. The index is aggregated to annual frequency for better visibility.

For each country, Index B measures the share of instability-related reporting relative to total news output over the preceding twelve months and is normalized by the average regional coverage of that dimension. This normalization expresses how strongly each country’s attention to a given instability dimension deviates from the regional norm, adjusted for differences in overall media activity.

The resulting index averages these normalized measures across all countries, assigning equal weight to each. The figure displays these averages separately for each dimension, with bars stacked by country, illustrating how regional attention to political instability has evolved and which countries have driven major shifts. Bars above one indicate periods of above-average regional attention, while bars below one mark calmer phases of reduced media salience.

Across dimensions, political violence shows relatively stable development over time and across countries. The region experienced some more violent years until the mid 1990s, a period of elevated violence after 2010, and then an increase in violence since 2019, which stayed elevated until recently, where a decline can be observed in recent years. Mass civil protests and instability of the regime experienced strong peaks in the 1990s and have decreased significantly since then. A few countries have mostly driven some recent upticks in the instability of the regime. Similarly, political instability within the regime shows rather low volatility with elevated values until the mid 1990s, but on a rather low level.


\begin{figure}[!t]
\centering
\caption{Average Intensity of Political Instability by Country (Index B)}
\includegraphics[width=\textwidth]{Figures/topic_intensity_by_year.pdf}
\label{fig:country_count_by_topic_stacked_avg}
\begin{minipage}{\textwidth}
 % Adds a bit of vertical space between the figure and the notes
    {\footnotesize{\textit{Notes:} This figure shows a composite index based on an equally weighted average of each country's Index B for the four political instability dimensions: political violence, mass civil protest, instability within the regime, and instability of the regime. It shows the contributions of the single countries as stacked bars.}}
\end{minipage}
\end{figure}



\subsection{Importance by Dimension}\label{sec:importance_by_dimension}


Figure \ref{fig:country_count_by_country} illustrates the evolution of media attention to political instability across West and Central African countries, disaggregated by four dimensions: political violence, mass civil protest, instability within the regime, and instability of the regime. Each bar shows a country’s annual coverage intensity for the four dimensions, normalized by the regional mean of that dimension, so that a value of one represents average regional attention. Bars above one indicate years when instability-related reporting in that country exceeded the regional norm, while lower bars denote calmer periods with below-average attention. 

The four dimensions of political instability -- political violence, mass civil protest, instability within the regime, and instability of the regime -- display distinct temporal and cross-country patterns. Political violence is the most persistent and geographically widespread component of instability, maintaining relatively stable but consistently elevated levels across most countries. Its prominence in the Central African Republic, Mali, Niger, and Chad reflects long-standing conflict dynamics and recurrent insurgent activity. 

Mass civil protest, by contrast, shows sharp spikes in specific years and countries, particularly during the early 1990s and mid-2010s, coinciding with waves of democratization and more recent episodes of social unrest. However, there is a clear decline in most countries afterward. This pattern suggests that protest-related instability tends to occur in short-lived waves rather than as a permanent feature of the political landscape.

Instability within the regime evolves more moderately but remains elevated throughout the observation period, indicating ongoing political contestation, factionalism, and leadership reshuffles within ruling elites -- most visibly in countries like Chad, Cameroon, and the Democratic Republic of the Congo. Instability of the regime, capturing overt challenges to government continuity such as coups or leadership overthrows, exhibits pronounced fluctuations: high in the 1990s, subdued in the 2000s, and rising again after 2020 with renewed coup activity in Mali, Niger, and Burkina Faso.
Instability within the regime and instability of the regime often rise simultaneously, indicating that elite fragmentation and regime threats frequently emerge alongside broader societal unrest and conflict. These dimensions act as amplifiers of instability during crises but fade when countries return to calmer conditions.

When compared across dimensions, the figure shows that political violence and regime instability tend to move together, with violent episodes often accompanying or following attempts to seize or retain power. Mass protest behaves less synchronously, frequently preceding institutional crises or occurring independently in relatively stable regimes. Political crises that produce widespread violence or protests often coincide with internal rifts within the ruling elite or direct threats to regime survival.




\begin{figure}[!t]
\centering
\caption{Average Intensity of Political Instability by Dimensions (Index B)}
\includegraphics[width=0.9\textwidth]{Figures/country_intensity_by_year.pdf}
\label{fig:country_count_by_country}
\begin{minipage}{\textwidth}
 % Adds a bit of vertical space between the figure and the notes
    {\footnotesize{\textit{Notes:} This figure shows a composite index based on an equally weighted average of each dimension's Index B for each country. It shows the contributions of the four political instability dimensions as stacked bars: political violence, mass civil protest, instability within the regime, and instability of the regime.}}
\end{minipage}
\end{figure}

\section{Empirical Analysis}\label{sec:empirical_analysis}

\subsection{Contagion Analysis}

As political instability might be, in some dimensions, a local phenomenon, there are reasons why it might be contagious and spread across countries. To assess the impact of political instability, it is relevant to question by how much political instability in one country is influenced by the degree of political instability in another country.

Giving a clear answer to this question is difficult and can be analyzed in different ways. A first approach is to look at the contemporaneous cross-country correlation of our political instability dimensions. Figure \ref{fig:correlation_heatmaps_annual} in the Appendix shows the cross-country correlation of the annually aggregated political instability indicators. We observe high contemporaneous correlations for some country combinations, which might sometimes be due to geographic, cultural, economic, or political connectedness. Looking at higher frequencies, such as months in figure \ref{fig:correlation_heatmaps_monthly} in the Appendix, results in lower correlation coefficients by construction. However, a high correlation still prevails for some country combinations, which may imply a certain level of connectedness. The lower correlation in higher frequencies might also mean that contagion can occur with a lag of several months, as information and actions take time to spread.

Incorporating the time structure of contagion, we can calculate the cross-correlation with different lags. Figures \ref{fig:contagion_network_geo_instability_of_regime}-\ref{fig:contagion_network_geo_political_violence} in the Appendix visualize the highest cross-correlation between each country within a maximum lag of 12 months. The figures show, for each dimension, which countries are most connected. We see an overall high correlation for political violence among the four dimensions, with higher connectedness among the core countries, supporting the geographical diffusion of violence. Similarly, mass civil protests have high correlations in some core countries that are often connected, while the peripheral countries show very low correlations. The dimensions of instability of and within the regime have lower correlations and seem, as a result, to be more affected by country-specific shocks.

% Running a VAR model and calculating the Impulse Responses between the political instability dimensions across countries gives an estimate of which dimension is more contagious to others, and for which countries this is most pronounced. Figure \ref{fig:irf_spillover_cross_indicator_NBS} and \ref{fig:irf_spillover_by_indicator_NBS} show them as cumulative Impulse responses over time.


\subsection{Macro Application: VAR Evidence on the Effects of Political Instability}\label{sec:maro_application}

%One of the paper’s most original contributions is its use of monthly VAR evidence, particularly with regard to the Democratic Republic of the Congo. The results clearly show differentiated macroeconomic responses depending on the type of political instability. Most strikingly, there is a counterintuitive finding that instability of the regime is associated with a temporary decline in policy rates. This is in contrast to the tightening observed after violence or mass protest shocks.

%This finding deserves substantially greater emphasis. It illustrates precisely why a multidimensional, high-frequency approach is important for real-time policy analysis, suggesting that central banks respond differently depending on the nature of political shocks. Similarly, the finding that economic activity (as measured by night-time lights) usually recovers within six to eight months provides a concrete timeframe for intervention that is relevant to policymakers.

In this section, we aim to evaluate the usability of our news-based political instability sentiment indicators for macroeconomic analysis. To do so, we examine the macroeconomic effects of an exogenous innovation to political instability by estimating a structural VAR model following \cite{caldara2022}. The reduced‐form VAR of order \(p\) is written as 
\begin{equation}
\mathbf{x}_t = \mathbf{c} + \sum_{\ell=1}^p A_\ell\,\mathbf{x}_{t-\ell} + \mathbf{u}_t
\end{equation}
where $\mathbf{x}_t$ is an $n \times 1$ vector of endogenous variables, $\mathbf{c}$ is an $n \times 1$ vector of constants, $A_\ell$ are $n \times n$ coefficient matrices for lags $\ell = 1, \ldots, p$, and $\mathbf{u}_t \sim \mathcal{N}(0, \Sigma)$ is an $n \times 1$ vector of reduced-form residuals with covariance matrix $\Sigma$. In this implementation, exogenous variables are treated as additional regressors in the system. 

The Bayesian estimation employs diffuse (non-informative) priors on both the coefficients and the error covariance. The prior on $\Sigma$ is specified with zero prior degrees of freedom, making the posterior dependent entirely on the data likelihood. Similarly, the prior precision on the stacked coefficient vector is zero, such that the posterior variance-covariance of the coefficients is proportional to $\Sigma \otimes (X'X)^{-1}$, where $X$ is the design matrix of regressors. We estimate the VAR models using Bayesian techniques with 20,000 draws from the posterior distribution of the structural parameters, employing a burn-in period of the first 4,000 draws.

Structural identification is achieved through a Cholesky decomposition of the covariance matrix of the VAR reduced-form residuals, following the literature showing that news sentiment reacts contemporaneously only to its own shock (see e.g. \cite{caldara2022, gambetti2023bad, LARSEN2019203}). The reduced-form residuals are orthogonalized via Cholesky decomposition,
\begin{equation}
\Sigma = LL',
\end{equation}
yielding structural shocks $\boldsymbol{\epsilon}_t = L^{-1}\mathbf{u}_t$ that are uncorrelated and have unit variance. We order the variables in the system according to the speed of their transmission, with the NBS index first, followed by the exchange rate, commodity price growth, consumer price inflation, the policy rate, and nighttime light as a proxy for economic activity, enabling impulse response analysis.


The dataset consists of: (i) the log of the NBS index; (ii) the log change in the exchange rate (LCU/USD); (iii) commodity price growth; (iv) annual consumer price inflation; (v) the policy rate of the central bank; (vi) real gross domestic product (GDP) growth; or nighttime light data as a high-frequency proxy of GDP (vii). Data is collected in annual, quarterly, and monthly frequencies, if available. See Table \ref{tab:macro_summary} for further details.\footnote{As we assume that commodity prices are relevant for the analyzed country, we tested the version with commodity price growth as an exogeneous variable in the system, since countries are most likely price takers. Different versions using the commodity price variable as endogenous and not using the variable are tested. The results stay the same. See \ref{fig:IRF_Democratic_Republic_of_the_Congo_GPRBASELINE_endo_com} and \ref{fig:IRF_Democratic_Republic_of_the_Congo_GPRBASELINE_no_com} in the Appendix.} %GDP. %In the CFA (African Financial Community) zone countries with fixed exchange rates and centralized monetary policy, we include the exchange rate and interest rate as exogeneous variables in the VAR specification, as these variables are exogenous to domestic shocks and do not respond to country-specific political or macroeconomic developments.

\begin{table}[h]
   \caption{Summary Macroeconomic Variables}
   \label{tab:macro_summary}
   \centering
   \scriptsize
   \resizebox{\linewidth}{!}{%
      \begin{tabular}{>{\raggedright\arraybackslash}p{3cm} >{\raggedright\arraybackslash}p{5cm} l l l l}
         \hline \addlinespace[0.3em]
         \hline
         Indicator & Unit & Frequency & Source\\ 
         \addlinespace[0.3em]
         \hline
         News-Based Sentiment & Log, normalized, 1990-M1 -- 2024-M12 & Monthly &  Own calculations \\
         Exchange Rate & Domestic currency per US Dollar, period average, period-over-period percent change & Monthly & \cite{IMFDataPortal} \\
         Commodity Export Price Index & Individual commodities weighted by ratio of exports to GDP,  June 2012 = 100, rolling weights & Monthly & \cite{IMFDataPortal} \\
         Consumer Price Index (CPI) & Reference period (2010=100), all items, period average, Year-over-year (YOY) percent change & Monthly & \cite{IMFDataPortal} \\
         Policy Rate & Percent per annum, period average  & Monthly &  \cite{IMFDataPortal} \\
        %Real GDP & Constant prices, domestic currency, year-on-year percent change & Annual & \cite{WorldBank2025WDI} \\
        
        Nighttime Lights & Median intensity, index, deseasonalized, 2012-M1 -- 2024-M12 & Monthly & \cite{Najjar2024Nightlight}  \\ 
         \hline
         \hline
      \end{tabular}
   }%
   \vspace{0.5em} % Optional spacing between table and notes
 \begin{minipage}{15cm}
 {\footnotesize \textit{Notes:} This table shows the sources, availability, frequency, and unit of the macroeconomic variables used in the VAR analysis. The frequency column shows the highest available frequency.}
 \end{minipage}
\end{table}

% Check if results change with net export weighted commodity prices
% WAEMU: Benin, Burkina Faso, Côte d'Ivoire, Guinea-Bissau, Mali, Niger, Senegal, Togo GDP data is year-over-year in quarters
% Nigeria, Ghana, Senegal data from IMF is quarter-over-quarter


%The figures \ref{fig:IRF_Central_African_Republic_GPRBASELINE}-\ref{fig:IRF_Central_African_Republic_GPRBASELINE} show the effect of a one standard deviation shock to the NBS index on a specific variable as impulse responses. The black solid lines depict the median impulse response to a one-standard-deviation increase in the NBS index, and the blue shaded bands illustrate the 68 and 90 percent pointwise credible sets.

%Identifying the impact of higher frequency shocks, as we identify in news-based political instability sentiment, is rather difficult when using low-frequency macroeconomic data. As with most West and Central African countries, macroeconomic indicators exist only at an annual frequency. In Figure \ref{fig:IRF_Central_African_Republic_GPRBASELINE} we show the Impulse Responses of a one standard deviation shock to the four different NBS indicators on CPI and GDP growth using data for 1994-2024 for the Central African Republic. Whereas it is not possible to extract a clear signal from annual data for many countries, as the shock is often not strong and long enough, for the Central African Republic, which faced an elongated period of high political conflict, we can observe how a shock to political instability significantly increases CPI inflation and decreases GDP growth in the first year, with a weaker countermovement as the measures are in annual growth rates. As the conflict dominated the entire country, we can clearly see a strong synchronous behavior among the four dimensions.

% Annual data:
% Benin: 1993-2023, Burkina Faso: 1991-2024, Cameroon: 1994-2024, CAR: 1994-2024, Chad: 1994-2024, Côte d'Ivoire: 1990-2024, DRC: 1994-2023, E.Guinea: 1994-2022, Gabon: 1994-2024, Ghana: 1990-2023, Guinea-Bissau: 1991-2024, Mali: 1991-2024, Mauritania: 1991-2023, Niger: 1990-2024, Nigeria: 1990-2024, Senegal: 1990-2024


% \begin{figure}[!t]
% \centering
% \caption{Impulse Responses to Political Instability in the Central African Republic}
% \includegraphics[width=\textwidth]{Figures/IRF_Central_African_Republic_GPRBASELINE.png}
% \label{fig:IRF_Central_African_Republic_GPRBASELINE}
% \begin{minipage}{\textwidth}
%  % Adds a bit of vertical space between the figure and the notes
%     {\footnotesize{\textit{Notes:} This figure shows the Impulse Responses to a one standard deviation shock to the four political instability indicators for the Central African Republic for CPI and GDP growth in annual frequency.}}
% \end{minipage}
% \end{figure}


%As for the West African Economic and Monetary Union (WAEMU) countries, quarterly macroeconomic data exists, we perform the same analysis for quarterly estimates of annual CPI inflation and quarterly GDP growth rates using data from 2012 to 2024. Figure \ref{fig:IRF_Togo_GPRBASELINE} shows the Impulse Responses for Togo. The results for the other WAEMU countries can be found in the Appendix. The responses vary depending on the historical background of each country. However, many countries experience increases in CPI inflation and lower GDP growth in the subsequent quarters following a political instability shock.

% Quarterly data: 
% Benin: 2013.07.01-2023.10.01, Burkina Faso: 2013.07.01-2024.10.01, Côte d'Ivoire: 2013.07.01-2024.10.01, Ghana: 2015.01.01-2024.07.01, Guinea-Bissau: 2013.07.01-2024.07.01, Mali: 2013.07.01-2024.10.01, Niger: 2013.07.01-2024.10.01, Nigeria:2011.01.01-2024.04.01, Senegal: 2014.01.01-2024.10.01, Togo: 2013.07.01-2024.10.01


% \begin{figure}[!t]
% \centering
% \caption{Impulse Responses to Political Instability in Togo}
% \includegraphics[width=\textwidth]{Figures/IRF_Togo_GPRBASELINE.png}
% \label{fig:IRF_Togo_GPRBASELINE}
% \begin{minipage}{\textwidth}
%  % Adds a bit of vertical space between the figure and the notes
%     {\footnotesize{\textit{Notes:} This figure shows the Impulse Responses to a one standard deviation shock to the four political instability indicators for Togo for CPI and Nighttime Light in quarterly frequency.}}
% \end{minipage}
% \end{figure}

To evaluate the usability of our high-frequency indicator, we analyze in a case study for the Democratic Republic of the Congo the impact of a political instability shock on macroeconomic variables at a monthly frequency. Figure \ref{fig:IRF_Democratic_Republic_of_the_Congo_GPRBASELINE} shows the effects of a one standard deviation shock to the NBS index on the exchange rate against the US Dollar, the trade-weighted commodity export price index, the consumer price inflation rate, the policy rate, and nighttime light intensity as Impulse Responses. The black solid lines depict the median impulse response to a one-standard-deviation increase in the NBS index, and the blue shaded bands illustrate the 68 and 90 percent pointwise credible sets. To make historical references, we show in Figure \ref{fig:drc_nbs_normalized_events_2012_2024_prettylabels_notitle_v2} the normalized political instability indicators for the Democratic Republic of the Congo with dashed vertical lines highlighting political episodes related to the four dimensions.



\begin{figure}[!t]
\centering
\caption{Normalized News-based Political Instability Indicators for the Democratic Republic of the Congo}
\includegraphics[width=\textwidth]{Figures/drc_normalized_nbs_events.pdf}
\label{fig:drc_nbs_normalized_events_2012_2024_prettylabels_notitle_v2}
\begin{minipage}{\textwidth}
 % Adds a bit of vertical space between the figure and the notes
    {\footnotesize{\textit{Notes:} This figure shows the normalized news-based political instability indicators for the Democratic Republic of the Congo. The indicators are normalized with the maximum value equals 1 for better visibility. Dashed lines are datted political episodes described in the text.}}
\end{minipage}
\end{figure}

%The political instability increases significantly and remains elevated for about four months. High political instability is followed by a short-term devaluation of the local currency and an increase in commodity prices. Consumer price inflation increases over time and remains elevated for longer. This is more pronounced for politically motivated violence, while for mass civil protests and instability within the regime, this effect is rather negligible; for instability of the regime, this effect is rather reversed. In line with this, the central bank reacts with an increase in policy rates for those that lead to higher inflation, while the instability of the regime shock leads to a decrease in the policy rate. For all four political instability measures, we observe a short-lived decrease in economic activity, proxied in this case by monthly nighttime light data, which normalizes again after six to eight months.


% Monthly data:
%Burkina Faso, Cameroon, CAR, Chad, Côte d'Ivoire, DRC,Gabon, Ghana, Guinea-Bissau, Mali, Mauritania, Niger, Nigeria, Senegal, Togo: 2012.01.01-2024.05.01
%Benin: 2012.01.01-2023.11.01
%Equatorial Guinea: 2012.01.01-2022.10.01



\begin{figure}[!t]
\centering
\caption{Impulse Responses to Political Instability in the Democratic Republic of the Congo}
\includegraphics[width=\textwidth, trim=0 0 0 1cm, clip]{Figures/IRF_Democratic_Republic_of_the_Congo_GPRBASELINE.png}
\label{fig:IRF_Democratic_Republic_of_the_Congo_GPRBASELINE}
\begin{minipage}{\textwidth}
 % Adds a bit of vertical space between the figure and the notes
    {\footnotesize{\textit{Notes:} This figure shows the Impulse Responses to a one standard deviation shock to the four political instability indicators for the Democratic Republic of the Congo for CPI and GDP growth in monthly frequency.}}
\end{minipage}
\end{figure}

The impulse responses to a political violence shock in Figure \ref{fig:IRF_Democratic_Republic_of_the_Congo_GPRBASELINE} show an immediate jump in the violence indicator that fades within a few months, paired with a short-run exchange-rate depreciation, a modest increase consumer price inflation and interest rates, and a temporary decline in nighttime lights that recovers slowly. This pattern is most naturally interpreted as reflecting recurrent conflict escalations that raise perceived risk and disrupt local production and trade routes, without usually translating into a nationwide institutional breakdown. % immediate supply disruption

In Figure \ref{fig:drc_nbs_normalized_events_2012_2024_prettylabels_notitle_v2}, major violence peaks coincide with concrete escalation moments -- most clearly late 2012 (M23/Goma) and mid-2022 (renewed M23 advances). These episodes plausibly map into the IRF mechanism as follows: the early exchange-rate depreciation reflects a rapid increase in perceived risk and precautionary FX demand; the modest CPI increase is consistent with temporary supply-chain and market-access disruptions; and the temporary reduction of nighttime lights is consistent with damage, displacement, or reduced activity in affected provinces that recover as acute risk of violence diminishes.

The impulse responses to a mass civil protest shock in Figure \ref{fig:IRF_Democratic_Republic_of_the_Congo_GPRBASELINE} display a sharp, short-lived spike in the protest indicator followed by immediate exchange-rate depreciation, rather weak to no effect on consumer price inflation with an insignificant increase of interest rates, and nighttime lights fall before gradually normalizing. This configuration matches the historical role of protest episodes as acute disruptions concentrated in economic hubs that can quickly increases perceived risk and FX pressure with potential to threaten political credibility in the longer term. % immediate risk increase

In Figure \ref{fig:drc_nbs_normalized_events_2012_2024_prettylabels_notitle_v2}, the largest protest clusters align with the 2015–2016 and late-2016 mobilizations tied to electoral delays and constitutional tensions, and with further waves in late 2017/early 2018 and the 2018–2019 transition period. The IRFs are consistent with these episodes operating through a fast uncertainty channel that shifts portfolios toward foreign currency, producing the early exchange-rate movement. The decline in nighttime lights then reflects the broader (though temporary) slowdown in activity during large, nationally salient protest waves.

The impulse responses to instability within the regime in Figure \ref{fig:IRF_Democratic_Republic_of_the_Congo_GPRBASELINE} show that the indicator spikes and decays, with exchange-rate pressure lasting longer, consumer price inflation weakly positive, no clear signal to interest rates, and nighttime lights fall mildly but more persistently. This pattern matches episodes where the political system continues functioning but policy-making becomes frictional -- delaying implementation, weakening credibility, and reducing investment -- rather than causing immediate nationwide shutdowns. %  medium term credibility loss
In Figure \ref{fig:drc_nbs_normalized_events_2012_2024_prettylabels_notitle_v2}, within-regime instability is elevated around the end-2016 electoral deadline crisis and remains salient through the early transition period in 2019, with another notable episode around late 2020 during coalition rupture and reconfiguration. These periods plausibly correspond to the IRF mechanism: uncertainty is not concentrated in a single day of disruption (as in protests) but in sustained bargaining and institutional conflict, which is why the exchange rate and nighttime lights react more slowly and persistently, and CPI drifts upward rather than jumping sharply.

The impulse responses to instability of the regime in Figure \ref{fig:IRF_Democratic_Republic_of_the_Congo_GPRBASELINE} are distinct to the others. The exchange rate depreciates only weakly and consumer price inflation even reacts negatively with interest rates falling persistently, and nighttime lights experience the most persistent decline. This pattern matches major political crises that make people doubt whether the country’s basic political rules will hold. In those moments, households and firms tend to pull back, cutting spending and delaying investment, so the economy slows mainly because confidence falls and demand weakens, not because supply chains are disrupted. In Figure \ref{fig:drc_nbs_normalized_events_2012_2024_prettylabels_notitle_v2}, regime-instability spikes concentrate around the late-2016 constitutional-electoral impasse and the 2018–2019 transition, with renewed stress around the 2023 election period. The IRFs map naturally onto these episodes: the initial exchange-rate depreciation captures immediate risk repricing; the CPI decline is consistent with households and firms postponing spending and investment, which can dominate supply pressures in a broad confidence shock; falling interest rates are consistent with weak credit demand and/or accommodation; and the persistent nighttime-lights decline reflects sustained reductions in real activity typical of system-wide uncertainty rather than localized disruption. %demand restraint

This case study shows how the multidimensional news-based political instability measures are able to separate distinct political instability shocks with different implications for the economy. Because the indicators can be constructed at a higher frequency, they sharpen identification in dynamic models and reveal heterogeneous transmission: violence shocks resemble localized risk and supply disruptions with limited persistence; protest shocks generate short term risk increases; within-regime instability maps to slower, more persistent uncertainty effects; and regime instability behaves like a broad confidence shock with demand restraint. Depending on the shock, different policy implications arise.

%For robustness, we tested various specifications, including other frequencies, other output measures, using variables as endogenous, exogenous variables, or exclude them completely, change the order of the variables in the system, and using different lag in the model. We draw similar conclusions from those robustness checks.






%\subsection{Micro Application: Political Instability and Firm-Level Effects}\label{sec:micro_application}



\section{Conclusion}\label{sec:conclusion}

This paper proposes a new framework for measuring political instability in West and Central African countries using high-frequency news reports. Building on the multidimensional conception of political instability and recent advances in news-based indicators, we construct monthly indices that separately track politically motivated violence, mass civil protest, instability within the political regime, and instability of the political regime. Our approach applies transparent keyword-based search strategies to generate indicators that are timely, comparable across countries, and suitable for econometric applications that require high-frequency inputs.

The empirical analysis shows that these four dimensions capture distinct yet interrelated facets of political instability and carry economically relevant information. They spike around different historical events, differ in persistence, and display cross-border as well as cross-dimension contagion. Political violence appears as a relatively persistent background risk, whereas mass civil protest and instability within the regime and of the regime are more episodic, often clustered around coups, contested elections, or constitutional reforms. %In country-specific VAR models, shocks to the NBS indices are associated with higher inflation and weaker real economic activity, with heterogeneous effects across countries and dimensions.
In a case study of the Democratic Republic of the Congo, embedding our indicators in a VAR model using high-frequency data, we find that our four political instability shocks affect macroeconomic variables in different ways: while political violence shocks resemble localized risk and supply disruptions; mass civil protests create short-term risk increases; within-regime instability maps slower, more persistent uncertainty; and regime instability behaves like a broad confidence shock with demand restraint.

%The firm-level analysis similarly indicates that periods of heightened political instability are linked to lower sales and employment, highlighting that political risks transmit to the microeconomic level and shape private sector performance.

At the same time, our framework has important limitations. First, the indicators are based on major English-language news outlets and therefore reflect the salience of political instability in international media rather than a complete census of events on the ground. Coverage intensity differs across countries and over time, and media biases linked to geographic, commercial, or political priorities may systematically influence the signal extracted from the data. Second, although the multidimensional structure allows us to distinguish between different forms of instability, it remains an approximation: some events are inherently ambiguous, and classification relies on keyword choices that may require further refinement. Finally, the macroeconomic analysis is constrained by the availability and quality of national accounts and monetary data, limiting the precision with which we can trace short-run transmission mechanisms. 


Despite these caveats, the proposed indicators provide a flexible and policy-relevant tool for monitoring and analyzing political instability in a region where traditional measures are often sparse, delayed, or overly aggregated. For policymakers and practitioners, they can complement existing event-based, expert-based, and survey-based datasets in country diagnostics, risk assessments, and early-warning exercises.


%Feedback:

%Gambia, Guinea, Liberia, Sierra Leone, Congo Republic, Rwanda, and Burundi could be included according to this definition.

% show an explicit sensitivity analysis quantifying how convergence with benchmark datasets varies between high- and low-coverage countries would be beneficial, as this would allow users to assess reliability across the sample.

%The authors should explicitly discuss whether Al Jazeera's involvement alters the composition of reporting, and whether this generates level shifts or changes in sensitivity over time. At the very least, this issue should be explicitly acknowledged as a potential source of a structural break.

%Firstly, the Cholesky ordering assumes that the news-based instability index is the fastest-moving variable in the system, preceding the exchange rate, inflation, the policy rate and output. This assumption is debatable. In particular, policy rates can often adjust more rapidly than inflation, so it is unclear whether the instability index is strictly exogenous to contemporaneous financial or policy variables. Alternative, more plausible orderings — especially those placing monetary policy earlier — should be more explicitly explored, at least as robustness checks.

%The cross-country contagion analysis provides compelling evidence of correlated instability across borders, particularly with regard to political violence. However, the analysis remains largely descriptive. The paper would be significantly strengthened by moving beyond correlation and testing specific transmission mechanisms. One promising approach would be to distinguish between geographic spillovers (e.g. shared borders or proximity) and institutional or political emulation channels (e.g. similarity in regime type or institutional fragility). Incorporating such a structure — either directly into the correlation analysis or through auxiliary regressions — would elevate this section from mapping connectedness to testing theories of diffusion. Similarly, the lag structure identified in the cross-correlation analysis could be used more explicitly to inform case-based narratives (e.g. violence propagating from Mali to Niger within a given timeframe), thereby integrating timing, geography and mechanism.

\spacingset{1}
\bibliographystyle{apalike}
\bibliography{literature_v1.5}

\clearpage

\appendix
\spacingset{1.84}

\setcounter{page}{1}
\setcounter{footnote}{0}

\renewcommand\thefigure{\thesection.\arabic{figure}} 
\renewcommand\thetable{\thesection.\arabic{table}}  
\setcounter{table}{0}
\setcounter{figure}{0}



%%%%%%%%%%%%%%%%%%%%%%%%%%%%%%%%%%%%%%%%%%%%%%%%%%%%%%%%%%%%%%%%%%%%%
\section{Additional Figures} \label{app:figures}
%%%%%%%%%%%%%%%%%%%%%%%%%%%%%%%%%%%%%%%%%%%%%%%%%%%%%%%%%%%%%%%%%%%%%
\FloatBarrier

\begin{figure}[h]
\centering
\caption{Map of Selected Countries}
\includegraphics[width=\textwidth]{Figures/africa_map.pdf}
\label{fig:africa_map}
\begin{minipage}{\textwidth}
    {\footnotesize{\textit{Notes:} The figure shows a map of Africa with the countries included in the study highlighted in red: Benin, Burkina Faso, Cameroon, Central African Republic, Chad, Côte d'Ivoire, Democratic Republic of the Congo, Equatorial Guinea, Gabon, Ghana, Guinea-Bissau, Mali, Mauritania, Niger, Nigeria, Senegal, and Togo.}}
\end{minipage}
\end{figure}

\begin{figure}[!t]
\centering
\caption{Number of Total Articles per Month}
\includegraphics[width=0.7\textwidth]{Figures/aggregate_normalization_reference_counts.pdf}
\label{fig:z_plot_all_cnt}
\begin{minipage}{\textwidth}
    {\footnotesize{\textit{Notes:} The figure shows the total number of monthly articles (excluding non-articles) for all countries included in the study. The black line shows monthly counts; the red line shows the twelve-month rolling mean.}}
\end{minipage}
\end{figure}

\begin{figure}[!t]
\centering
\caption{Comparison of Index A and Index B: Politically Motivated Violence}
\includegraphics[width=\textwidth]{Figures/nbs_ab_political_violence_by_country.pdf}
\label{fig:country_count_norm_AB_political_violence}
\begin{minipage}{\textwidth}
    {\footnotesize{\textit{Notes:} This figure compares the news-based political instability sentiment indicators (Index A and Index B) for the dimension of politically motivated violence across all sample countries.}}
\end{minipage}
\end{figure}

\begin{figure}[!t]
\centering
\caption{Comparison of Index A and Index B: Mass Civil Protest}
\includegraphics[width=\textwidth]{Figures/nbs_ab_mass_civil_protest_by_country.pdf}
\label{fig:country_count_norm_AB_mass_civil_protest}
\begin{minipage}{\textwidth}
    {\footnotesize{\textit{Notes:} This figure compares the news-based political instability sentiment indicators (Index A and Index B) for the dimension of mass civil protest across all sample countries.}}
\end{minipage}
\end{figure}

\begin{figure}[!t]
\centering
\caption{Comparison of Index A and Index B: Instability of Regime}
\includegraphics[width=\textwidth]{Figures/nbs_ab_instability_of_regime_by_country.pdf}
\label{fig:country_count_norm_AB_instability_of_regime}
\begin{minipage}{\textwidth}
    {\footnotesize{\textit{Notes:} This figure compares the news-based political instability sentiment indicators (Index A and Index B) for the dimension of instability of the regime across all sample countries.}}
\end{minipage}
\end{figure}

\begin{figure}[!t]
\centering
\caption{Comparison of Index A and Index B: Instability Within Regime}
\includegraphics[width=\textwidth]{Figures/nbs_ab_instability_within_regime_by_country.pdf}
\label{fig:country_count_norm_AB_instability_within_regime}
\begin{minipage}{\textwidth}
    {\footnotesize{\textit{Notes:} This figure compares the news-based political instability sentiment indicators (Index A and Index B) for the dimension of instability within the regime across all sample countries.}}
\end{minipage}
\end{figure}

\begin{figure}[!t]
\centering
\caption{Annual Cross-Country Correlation by Political Instability Dimension}
\includegraphics[width=\textwidth]{Figures/annual_correlation_heatmaps.pdf}
\label{fig:correlation_heatmaps_annual}
\begin{minipage}{\textwidth}
    {\footnotesize{\textit{Notes:} This figure shows annual cross-country correlation coefficients for each political instability dimension. Blue indicates negative correlations; red indicates positive correlations. Country codes: BFA = Burkina Faso, CMR = Cameroon, CAF = Central African Republic, TCD = Chad, CIV = Côte d'Ivoire, COD = Democratic Republic of the Congo, GNQ = Equatorial Guinea, GAB = Gabon, GHA = Ghana, GNB = Guinea-Bissau, MLI = Mali, MRT = Mauritania, NER = Niger, NGA = Nigeria, SEN = Senegal.}}
\end{minipage}
\end{figure}

\begin{figure}[!t]
\centering
\caption{Monthly Cross-Country Correlation by Political Instability Dimension}
\includegraphics[width=\textwidth]{Figures/monthly_correlation_heatmaps.pdf}
\label{fig:correlation_heatmaps_monthly}
\begin{minipage}{\textwidth}
    {\footnotesize{\textit{Notes:} This figure shows monthly cross-country correlation coefficients for each political instability dimension. Blue indicates negative correlations; red indicates positive correlations. Country codes as in Figure~\ref{fig:correlation_heatmaps_annual}.}}
\end{minipage}
\end{figure}


\begin{figure}[!t]
\centering
\caption{Geographical Contagion Network: Instability of Regime}
\includegraphics[width=\textwidth, trim=0cm 0cm 0cm 1.5cm, clip]{Figures/contagion_network_geo_instability_of_regime.pdf}
\label{fig:contagion_network_geo_instability_of_regime}
\begin{minipage}{\textwidth}
    {\footnotesize{\textit{Notes:} This figure displays the geographical contagion network for the regime instability dimension. Edge weights represent cross-correlation strength at optimal lag (up to 12 months). Country codes as in Figure~\ref{fig:correlation_heatmaps_annual}.}}
\end{minipage}
\end{figure}


\begin{figure}[!t]
\centering
\caption{Geographical Contagion Network: Instability Within Regime}
\includegraphics[width=\textwidth, trim=0cm 0cm 0cm 1.5cm, clip]{Figures/contagion_network_geo_instability_within_regime.pdf}
\label{fig:contagion_network_geo_instability_within_regime}
\begin{minipage}{\textwidth}
    {\footnotesize{\textit{Notes:} This figure displays the geographical contagion network for the within-regime instability dimension. Edge weights represent cross-correlation strength at optimal lag (up to 12 months). Country codes as in Figure~\ref{fig:correlation_heatmaps_annual}.}}
\end{minipage}
\end{figure}

\begin{figure}[!t]
\centering
\caption{Geographical Contagion Network: Mass Civil Protest}
\includegraphics[width=\textwidth, trim=0cm 0cm 0cm 1.5cm, clip]{Figures/contagion_network_geo_mass_civil_protest.pdf}
\label{fig:contagion_network_geo_mass_civil_protest}
\begin{minipage}{\textwidth}
    {\footnotesize{\textit{Notes:} This figure displays the geographical contagion network for the mass civil protest dimension. Edge weights represent cross-correlation strength at optimal lag (up to 12 months). Country codes as in Figure~\ref{fig:correlation_heatmaps_annual}.}}
\end{minipage}
\end{figure}

\begin{figure}[!t]
\centering
\caption{Geographical Contagion Network: Political Violence}
\includegraphics[width=\textwidth, trim=0cm 0cm 0cm 1.5cm, clip]{Figures/contagion_network_geo_political_violence.pdf}
\label{fig:contagion_network_geo_political_violence}
\begin{minipage}{\textwidth}
    {\footnotesize{\textit{Notes:} This figure displays the geographical contagion network for the political violence dimension. Edge weights represent cross-correlation strength at optimal lag (up to 12 months). Country codes as in Figure~\ref{fig:correlation_heatmaps_annual}.}}
\end{minipage}
\end{figure}


\begin{figure}[!t]
\centering
\caption{Impulse Responses to Political Instability in the Democratic Republic of the Congo}
\includegraphics[width=\textwidth, trim=0 0 0 1cm, clip]{Figures/IRF_Democratic_Republic_of_the_Congo_GPRBASELINE_endo_com.png}
\label{fig:IRF_Democratic_Republic_of_the_Congo_GPRBASELINE_endo_com}
\begin{minipage}{\textwidth}
 % Adds a bit of vertical space between the figure and the notes
    {\footnotesize{\textit{Notes:} This figure shows the Impulse Responses to a one standard deviation shock to the four political instability indicators for the Democratic Republic of the Congo for the change of the exchange rate, consumer price inflation, interest rates, and economic activity in monthly frequency using the commodity prices as endogeneous variables in the system.}}
\end{minipage}
\end{figure}

\begin{figure}[!t]
\centering
\caption{Impulse Responses to Political Instability in the Democratic Republic of the Congo}
\includegraphics[width=\textwidth, trim=0 0 0 1cm, clip]{Figures/IRF_Democratic_Republic_of_the_Congo_GPRBASELINE_no_com.png}
\label{fig:IRF_Democratic_Republic_of_the_Congo_GPRBASELINE_no_com}
\begin{minipage}{\textwidth}
 % Adds a bit of vertical space between the figure and the notes
    {\footnotesize{\textit{Notes:} This figure shows the Impulse Responses to a one standard deviation shock to the four political instability indicators for the Democratic Republic of the Congo for the change of the exchange rate, consumer price inflation, interest rates, and economic activity in monthly frequency using no commodity prices.}}
\end{minipage}
\end{figure}

%%%%%%%%%%%%%%%%%%%%%%%%%%%%%%
% force all pending floats to appear before moving on
\FloatBarrier
%%%%%%%%%%%%%%%%%%%%%%%%%%%%%%

%%%%%%%%%%%%%%%%%%%%%%%%%%%%%%%%%%%%%%%%%%%%%%%%%%%%%%%%%%%%%%%%%%%%%
\section{Additional Tables} \label{app:tables}
%%%%%%%%%%%%%%%%%%%%%%%%%%%%%%%%%%%%%%%%%%%%%%%%%%%%%%%%%%%%%%%%%%%%%

\begin{table}[h]
   \caption{Summary Statistics: Monthly Article Counts}
   \label{tab:summary_statistics}
   \centering
      \scriptsize
  \resizebox{0.7\linewidth}{!}{%
\begin{tabular}{lrrrr}
  \hline  \addlinespace[0.3em]
 \hline
Country & Mean & Median & p10 & p90 \\ 
  \addlinespace[0.3em]
 \hline
Nigeria & 445 & 448 & 123 & 706 \\ 
  Democratic Republic of the Congo & 193 & 162 & 43 & 366 \\ 
  Côte d'Ivoire & 132 & 86 & 32 & 219 \\ 
  Ghana & 92 & 79 & 27 & 151 \\ 
  Mali & 83 & 44 & 9 & 190 \\ 
  Senegal & 73 & 58 & 15 & 139 \\ 
  Cameroon & 58 & 47 & 16 & 107 \\ 
  Niger & 45 & 31 & 11 & 82 \\ 
  Burkina Faso & 40 & 23 & 5 & 85 \\ 
  Chad & 40 & 28 & 7 & 82 \\ 
  Central African Republic & 33 & 19 & 4 & 69 \\ 
  Togo & 27 & 15 & 5 & 49 \\ 
  Gabon & 25 & 17 & 7 & 41 \\ 
  Mauritania & 24 & 16 & 5 & 50 \\ 
  Guinea-Bissau & 22 & 10 & 2 & 63 \\ 
  Benin & 15 & 11 & 4 & 30 \\ 
  Equatorial Guinea & 15 & 8 & 1 & 32 \\ 
     \hline
  mean & 80 & 65 & 19 & 145 \\ 
median & 40 & 28 & 7 & 82 \\ 
  min & 15 & 8 & 1 & 30 \\ 
  max & 445 & 448 & 123 & 706 \\ 
   \hline
 \hline
\end{tabular}
  }%
     \vspace{0.5em} % Optional spacing between table and notes
 \begin{minipage}{0.7\linewidth}
 {\footnotesize \textit{Notes:} This table presents the summary statistics for the counted number of articles. The mean, median, 10th, and 90th percentiles of the number of articles are shown for each country, the mean, median, minimum, and maximum over all countries.}
 \end{minipage}
 \end{table}

\begin{table}[htbp]
\scriptsize
\centering
\caption{Mean Correlations across Countries}
\label{tab:mean_correlation_selected_events_ordered}
\begin{minipage}{0.7\linewidth}
\begin{tabular}{llrrrr}
\toprule
& & \multicolumn{4}{c}{NBS Dimensions} \\
\cmidrule(lr){3-6}
Event Type & Source & Violence & Protest & Within & Regime \\
\midrule
Political Violence &  & 1.00 & --- & --- & --- \\ 
  Mass Civil Protest &  & 0.57 & 1.00 & --- & --- \\ 
  Instability within Regime &  & 0.63 & 0.56 & 1.00 & --- \\ 
  Instability of Regime &  & 0.42 & 0.39 & 0.58 & 1.00 \\ 
  \midrule
  Civil Unrest & BARRET & 0.37 & 0.43 & 0.52 & 0.35 \\ 
  Internal Conflict & IRCG & 0.19 & 0.14 & 0.14 & 0.12 \\ 
  Coups & COUPS & 0.14 & 0.13 & 0.14 & 0.14 \\ 
  Annual Coups & CSP & 0.07 & 0.03 & 0.03 & 0.03 \\ 
  Civil War & IRCG & 0.05 & -0.06 & -0.05 & -0.08 \\ 
  Viol. & WHIV & 0.31 & 0.29 & 0.24 & 0.19 \\ 
  Viol. against Civilians & ACLED & 0.17 & 0.03 & 0.09 & 0.05 \\ 
  Ethnic Viol. & UCDP & 0.04 & -0.00 & 0.02 & 0.00 \\ 
  Battles & ACLED & 0.28 & 0.12 & 0.12 & 0.08 \\ 
  Explosions & ACLED & 0.14 & 0.03 & 0.04 & 0.01 \\ 
  Terrorism & IRCG & 0.05 & -0.06 & -0.05 & -0.08 \\ 
  Protest & WHIV & 0.11 & 0.24 & 0.16 & 0.13 \\ 
  Protests & ACLED & 0.16 & 0.16 & 0.13 & 0.16 \\ 
  Riots & ACLED & 0.12 & 0.11 & 0.08 & 0.07 \\ 
  Pol. Sanctions & WHIV & 0.10 & 0.19 & 0.15 & 0.10 \\ 
  Corruption Index & GALLUP & -0.03 & 0.05 & -0.04 & -0.02 \\ 
  Gov. Crises & CNTS & 0.02 & -0.01 & -0.00 & 0.03 \\ 
  Years in Office & DPI & 0.01 & 0.03 & -0.02 & -0.09 \\ 
  Legis. Frac. & DPI & -0.11 & -0.18 & -0.17 & -0.13 \\ 
  Democracy Index & CSP & -0.16 & -0.13 & -0.08 & 0.00 \\ 
  Gov. Effect. & WWGI & -0.06 & -0.06 & -0.07 & -0.09 \\ 
  Party in Power & DPI & -0.00 & 0.01 & -0.09 & -0.15 \\ 
  Maj. Party & DPI & -0.01 & 0.03 & -0.04 & -0.19 \\ 
  Pol. Stab. No Viol. & WWGI & -0.15 & 0.01 & 0.03 & -0.02 \\ 
  Reg. Quality & WWGI & -0.08 & 0.00 & -0.08 & -0.06 \\ 
  Rule of Law & WWGI & -0.00 & -0.03 & -0.03 & -0.08 \\ 
  Voice and Account. & WWGI & -0.05 & -0.06 & -0.03 & -0.08 \\ 
\bottomrule
\end{tabular}
\vspace{0.5em}
{\footnotesize \flushleft \textit{Notes:} This table reports mean pairwise correlations across countries between the four NBS dimensions Index B (Political Violence, Mass Civil Protest, Instability within Regime, and Instability of Regime) and a set of alternative political (in-)stability measures. Event Type denotes the name of the indicator, and Sources denotes the underlying dataset. The NBS Dimensions were aggregated to the same frequency as the alternative measure, and the time span was adjusted to the maximum common length available. Not all alternative measures exist for all countries.}
\end{minipage}
\end{table}


\begin{landscape}
\begin{table}[!t]
   \caption{Summary Statistics for All Dimensions: Monthly Article Counts}
   \label{tab:summary_statistics_all}
   \centering
   \resizebox{1.5\textwidth}{!}{%
 \begin{tabular}{lrrrrrrrrrrrrrrrrrrrrrrrr}
 \hline
 \multicolumn{1}{c}{} &  \multicolumn{4}{c}{All Articles} & \multicolumn{4}{c}{Political Instability} &  \multicolumn{4}{c}{Political Violence} &  \multicolumn{4}{c}{Mass Civil Protest} &  \multicolumn{4}{c}{Instability Within Regime} &  \multicolumn{4}{c}{Instability Of Regime} \\
 \hline
 Country & Mean & Median & p10 & p90& Mean & Median & p10 & p90 & Mean & Median & p10 & p90 & Mean & Median & p10 & p90 & Mean & Median & p10 & p90 & Mean & Median & p10 & p90 \\
 \hline
  \addlinespace[0.3em]
 \hline
Benin & 15 & 11 & 4 & 30 & 7 & 3 & 0 & 16 & 3 & 1 & 0 & 9 & 1 & 0 & 0 & 1 & 3 & 1 & 0 & 6 & 0 & 0 & 0 & 1 \\ 
  Burkina Faso & 40 & 23 & 5 & 85 & 34 & 10 & 0 & 83 & 26 & 6 & 0 & 74 & 3 & 0 & 0 & 6 & 6 & 2 & 0 & 12 & 2 & 0 & 0 & 4 \\ 
  Cameroon & 58 & 47 & 16 & 107 & 27 & 10 & 1 & 71 & 21 & 7 & 0 & 60 & 2 & 0 & 0 & 6 & 4 & 2 & 0 & 10 & 0 & 0 & 0 & 1 \\ 
  Central African Republic & 33 & 19 & 4 & 69 & 43 & 15 & 1 & 103 & 32 & 10 & 0 & 87 & 1 & 0 & 0 & 4 & 8 & 3 & 0 & 18 & 2 & 0 & 0 & 5 \\ 
  Chad & 40 & 28 & 7 & 82 & 32 & 17 & 3 & 71 & 24 & 13 & 2 & 55 & 2 & 0 & 0 & 4 & 6 & 3 & 0 & 13 & 1 & 0 & 0 & 3 \\ 
  Côte d'Ivoire & 132 & 86 & 32 & 219 & 109 & 28 & 7 & 200 & 75 & 17 & 4 & 135 & 10 & 1 & 0 & 16 & 25 & 8 & 1 & 48 & 5 & 1 & 0 & 14 \\ 
  Democratic Republic of the Congo & 193 & 162 & 43 & 366 & 152 & 108 & 23 & 339 & 115 & 83 & 16 & 259 & 7 & 3 & 0 & 17 & 28 & 13 & 3 & 67 & 6 & 2 & 0 & 12 \\ 
  Equatorial Guinea & 15 & 8 & 1 & 32 & 4 & 1 & 0 & 12 & 2 & 0 & 0 & 6 & 0 & 0 & 0 & 1 & 2 & 0 & 0 & 4 & 0 & 0 & 0 & 1 \\ 
  Gabon & 25 & 17 & 7 & 41 & 12 & 4 & 0 & 20 & 6 & 2 & 0 & 10 & 1 & 0 & 0 & 2 & 5 & 1 & 0 & 8 & 1 & 0 & 0 & 2 \\ 
  Ghana & 92 & 79 & 27 & 151 & 17 & 9 & 1 & 34 & 9 & 5 & 0 & 21 & 1 & 0 & 0 & 2 & 6 & 3 & 0 & 13 & 1 & 0 & 0 & 3 \\ 
  Guinea-Bissau & 22 & 10 & 2 & 63 & 20 & 5 & 0 & 44 & 13 & 3 & 0 & 29 & 1 & 0 & 0 & 3 & 6 & 2 & 0 & 13 & 1 & 0 & 0 & 3 \\ 
  Mali & 83 & 44 & 9 & 190 & 78 & 21 & 2 & 210 & 61 & 15 & 1 & 166 & 4 & 1 & 0 & 10 & 12 & 3 & 0 & 31 & 3 & 0 & 0 & 7 \\ 
  Mauritania & 24 & 16 & 5 & 50 & 15 & 6 & 0 & 37 & 9 & 3 & 0 & 25 & 1 & 0 & 0 & 3 & 4 & 1 & 0 & 9 & 1 & 0 & 0 & 2 \\ 
  Niger & 45 & 31 & 11 & 82 & 35 & 18 & 3 & 72 & 25 & 13 & 2 & 56 & 2 & 1 & 0 & 5 & 7 & 2 & 0 & 15 & 2 & 0 & 0 & 4 \\ 
  Nigeria & 445 & 448 & 123 & 706 & 240 & 192 & 47 & 456 & 177 & 138 & 28 & 376 & 18 & 10 & 1 & 44 & 47 & 30 & 10 & 91 & 7 & 5 & 1 & 15 \\ 
  Senegal & 73 & 58 & 15 & 139 & 29 & 17 & 4 & 57 & 17 & 10 & 2 & 34 & 3 & 1 & 0 & 5 & 9 & 4 & 0 & 18 & 1 & 0 & 0 & 4 \\ 
  Togo & 27 & 15 & 5 & 49 & 19 & 5 & 0 & 49 & 10 & 2 & 0 & 24 & 2 & 0 & 0 & 6 & 6 & 1 & 0 & 14 & 1 & 0 & 0 & 3 \\ 
  mean & 80 & 65 & 19 & 145 & 51 & 28 & 5 & 110 & 37 & 19 & 3 & 84 & 3 & 1 & 0 & 8 & 11 & 5 & 1 & 23 & 2 & 0 & 0 & 5 \\ 
   \hline
 \hline
median & 40 & 28 & 7 & 82 & 29 & 10 & 1 & 71 & 21 & 7 & 0 & 55 & 2 & 0 & 0 & 5 & 6 & 2 & 0 & 13 & 1 & 0 & 0 & 3 \\ 
  min & 15 & 8 & 1 & 30 & 4 & 1 & 0 & 12 & 2 & 0 & 0 & 6 & 0 & 0 & 0 & 1 & 2 & 0 & 0 & 4 & 0 & 0 & 0 & 1 \\ 
  max & 445 & 448 & 123 & 706 & 240 & 192 & 47 & 456 & 177 & 138 & 28 & 376 & 18 & 10 & 1 & 44 & 47 & 30 & 10 & 91 & 7 & 5 & 1 & 15 \\ 
   \hline
 \hline
\end{tabular}
  }%
     \vspace{0.5em}
 \begin{minipage}{0.7\linewidth}
 {\footnotesize \textit{Notes:} This table presents summary statistics for the monthly count of articles related to political instability, disaggregated by dimension.}
 \end{minipage}
 \end{table}
\end{landscape}




\begin{table}[h!]
\caption{Search Criteria}
\label{tab:search_terms}
\centering
\scriptsize
\begin{tabular}{|p{0.1\linewidth}|p{0.8\linewidth}|}
\hline
\textbf{Type} & \textbf{Content} \\ \hline
\textbf{Exclude Tags} & Abstracts, Advertorials/Sponsored Content, AI-Generated Content, Audio-visual Links, Calendar and Events, Chronology, Columns, Commentaries/Opinions, Corporate Digest, Country Profiles, Economic Predictions/Forecasts, Editorials, Headline Listings, Headline-Only Content, Images, Interviews, Letters, News Digest, Obituaries, People Profiles, Personal Announcements, Press Releases, Prospectuses, Rankings, Reviews, Routine General News, Routine Market, Statistics, Surveys/Polls, Tables, Transcripts, Commodity/Financial Market News, Weather, Sports, Books, Games, Movies, Theater, International Relations, Space, Royal Families, Sciences/Humanities, Real Estate/Property, Animal Attacks, Arts/Entertainment, Awards, Fire/Rescue Services, Health, Global World Issues, Economic Performance/Indicators, Trade/External Payments, Eurozone Data/Policies, Economic News, Corporate News \\ \hline
\textbf{Include Tags} & Violence based search: \begin{itemize}
    \item Domestic Politics, Public Security, Civil Unrest, Armed Forces, Terrorism
     \end{itemize}
     Politics based search:
      \begin{itemize}
    \item Domestic Politics, Public Security, Elections, impeachment, Political Appointments/Terminations
\end{itemize}
\\ \hline
\textbf{Exclude Strings} & \cite{caldara2022}: \begin{itemize}
 \item movie*, film*, museum*, anniversar*, obituar*, memorial*, arts, book, books, memoir*, ``price war'', game, story, history, veteran*, tribute*, sport, music, racing, cancer, ``real estate'', mafia, trial, tax
 \end{itemize}
 \cite{BARRETT2022}: \begin{itemize}
 \item ``vote of protest'', ``protest vote'', protestant*, do not use ``demonstration'' 
\end{itemize} \\ \hline
\textbf{Include Strings} & 

Politically motivated Violence: 
   
Guerilla Warfare \newline
   ((guerrilla OR guerilla OR militant OR rebel* OR ``armed resistance'' OR militia OR paramilitary OR Islamist* OR terrorist* OR jihadist OR separatist OR ``Amabazonia Defence Forces'' OR ADF OR ``Boko Haram'' OR CPDM OR Seleka OR ``Anti-balaka'' OR UPC OR FPRC OR JNIM OR ``Islamic State'' OR ISGS OR ``Forces Nouvelles'' OR GMDF OR MNLA OR MFDC OR AQIM OR ``Islamic Salvation Front'' OR FIS) 
   AND (sabotage OR insurgen* OR ambush OR overthrow OR coup OR rebellion OR ``regime change'' OR ``irregular warfare'' OR ``unconventional warfare'' OR battle OR combat OR terror OR attack* OR kill* OR murder* OR shot* OR raid OR die* OR ``under fire'' OR losses)) \newline 

Revolutions \newline
   (revolution* OR ``armed uprising'' OR rebellion OR ``armed insurrection'' OR mutiny* OR revolt OR overthrow OR ``independence movement'' OR secession OR ``freedom fight'' OR ``power shift'' OR ``political upheaval'' OR ``coup plot'') NEAR10 (anti-government OR government OR ``democratic movement'' OR liberation OR opposition OR dissident OR transformation OR ``political change'' OR ``regime change'' OR (end w/3 military) OR (end w/3 dictator*) OR (end w/3 authoritarian*))/F50/ \newline

Civil War \newline
   (``civil war'' OR ``internal conflict'' OR ``civil strife'' OR ``domestic conflict'' OR ``armed rebellion'' OR ``internal warfare'' OR ``national crisis'' OR ``armed conflict'' OR ``government crackdown'' OR ``conflict escalation'')/F50/ AND (``mass casualties'' OR ``military engagement'' OR separatist OR government OR kill* OR die* OR purge* OR death OR murder* OR violen*)/F50/ \newline

      \\ \hline
   \end{tabular}
\end{table}


\begin{table}[h]
   \caption*{Search Criteria (continued)}
   \centering
   \scriptsize
   \begin{tabular}{|p{0.1\linewidth}|p{0.8\linewidth}|}
   \hline
   \textbf{Type} & \textbf{Content} \\ \hline
   \textbf{Include Strings} & 
Internal Conflicts \newline
   (conflict OR violen* OR war OR rebellion OR crisis OR assault OR murder* OR kill* OR terror* OR attack* OR ambush OR insurgen*) 
   AND (politic* OR government OR ``government official'' OR ``public official'' OR ``political leader'' OR representative OR congressman OR ``cabinet member'' OR minister OR governor OR senator OR lawmaker OR mayor OR ``member of parliament'' OR president) \newline 


Ethnic Tensions \newline
   (``Ethnic tensions'' OR ``ethnic violence'' OR ``Racial tensions'' OR ``Ethnic conflict'' OR ``Racial conflict'' OR ``Nationality divisions'' OR ``Interethnic violence'' OR ``Racial unrest'' OR ``Ethnic unrest'' OR ``Cultural tensions'' OR ``Identity conflict'' OR ``Ethnic strife'' OR ``Racial strife'' OR ``Communal violence'' OR ``Sectarian tensions'' OR ``Ethnic discrimination'' OR ``Racial discrimination'' OR ``National identity crisis'' OR ``Ethnic polarization'' OR (ethnic OR racial) w/3 (violence OR conflict))/F50/ \newline 

Riots \newline
   (riot OR ``civil unrest'' OR ``public disturbance'' OR ``violent protest'' OR ``street violence'' OR ``mob violence'' OR uprising OR ``urban unrest'' OR ``crowd violence'' OR ``mass riot'' OR ``social unrest'' OR ``civil disorder'' OR ``riotous behavior'' OR ``civil disturbance'' OR ``violent clash'' OR ``public disorder'' OR ``riot outbreak'' OR ``mass disturbance'' OR ``public riot'')/F50/ \newline 

Demonstrations \newline
   (``Protest'' OR ``Demonstration'' OR ``Sit-in'' OR ``Peaceful protest'' OR ``Public gathering'' OR ``Protest march'' OR ``Protest rally'' OR ``Peaceful demonstration'' OR ``Mass protest'' OR ``Public demonstration'' OR ``Activist gathering'' OR ``Protest event'' OR ``Political protest'' OR ``Anti-government protest'' OR ``Civic protest'' OR ``Social movement'' OR ``Public dissent'' OR ``Protest action'') \newline 

Strikes \newline
   (``Work stoppage'' OR ``Labor strike'' OR ``Industrial action'' OR Walkout OR ``General strike'' OR ``Employee strike'' OR ``Workers' strike'' OR ``Job action'' OR ``Strike action'' OR ``Work protest'' OR ``Union strike'' OR ``Labor protest'' OR ``Trade union strike'' OR ``Collective action'' OR ``Sit-down strike'' OR ``Strike movement'' OR ``Workplace strike'' OR ``Strike rally'' OR ``Workforce protest'' OR ``union protest'' OR ``wage protest'' OR ``employment protest'') \newline 

Instability within the political regime \newline
   (democracy OR elect* OR vote* OR multiparty OR ``political alliance'' OR suffrage OR ``political parties'' OR ``government stability'' OR turnout OR multiparty OR ``political competition'' OR political w/3 fragmentation OR ballot OR ``coalition government'' OR ``electoral alliances'' OR ``electoral diversity'' OR contender* OR ``political alliance*'' OR ``political diversity'' OR ``political pluralism'' OR democratic w/3 (election OR vote OR right OR reform OR institution))/F50/ \newline

 Instability of the political regime \newline
    ``new constitution'' OR ``cabinet change'' OR coups d'etat OR ``constitutional change'' OR anarchy OR ``constitutional reform'' OR ``regime change'' OR ``electoral reforms'' OR ``election disputes'' OR ``constitutional referendum''  OR (new OR change or transition OR amend* OR reform OR overhaul OR restructure) w/3 (political OR constitutional OR leader* OR electoral OR systemic OR policy OR state OR government) \\
 \hline
\end{tabular}
   \vspace{0.5em}
   \begin{minipage}{\textwidth}
   {\footnotesize \flushleft \textit{Notes:} This table summarizes the search terms used to construct the news-based political instability indices. Search operators include ``and'', ``or'', ``not'', ``near'', ``atleast'', ``within'', truncation sign ``*'', first x words, word count, and question mark.}
   \end{minipage}
\end{table}


%%%%%%%%%%%%%%%%%%%%%%%%%%%%%%%%%%%%%%%%%%%%%%%%%%%%%%%%%%%%%%%%%%%%%
% IMPROVED POLITICAL EVENTS TABLE WITH PRIMARY DIMENSION COLUMN
%%%%%%%%%%%%%%%%%%%%%%%%%%%%%%%%%%%%%%%%%%%%%%%%%%%%%%%%%%%%%%%%%%%%%
\newcolumntype{L}[1]{>{\raggedright\arraybackslash}p{#1}}
\renewcommand{\arraystretch}{0.85}
\setlength{\tabcolsep}{3pt}
\begin{table}[htbp]
\centering
\scriptsize
\caption{Major Political Instability Events and NBS Index Response, 1990--2025}
\begin{tabular}{p{2.5cm} p{1.5cm} p{7.5cm} L{3cm}}
\toprule
\textbf{Country} & \textbf{Date} & \textbf{Key Event} & \textbf{Primary NBS Dimension} \\
\midrule
Benin & 1990--1991 & Transition to multiparty democracy; National Conference & Within-Regime \\
      & 2019 & Protests over exclusionary parliamentary elections & Protest \\
\midrule
Burkina Faso & 2014--10 & Popular uprising ousts President Compaoré & Protest, Regime \\
              & 2015--09 & Failed coup attempt against interim government & Violence, Regime \\
              & 2022--01 & Military coup removes President Kaboré & Regime \\
              & 2022--09 & Second coup; Captain Traoré assumes power & Regime \\
\midrule
Cameroon & 2016--10 & Anglophone protests escalate into armed conflict & Violence, Protest \\
         & 2018--10 & Contested presidential election amid separatist violence & Violence, Within \\
\midrule
Central African & 2003--03 & Bozizé coup ousts President Patassé & Violence, Regime \\
Republic & 2013--03 & Séléka rebellion captures Bangui; civil war begins & Violence, Regime \\
         & 2019--02 & Khartoum peace accord signed & Violence \\
\midrule
Chad & 2006--04 & Rebel assault on N'Djamena repelled & Violence \\
     & 2021--04 & President Déby killed; military council seizes power & Violence, Regime \\
\midrule
Côte d'Ivoire & 2002--09 & Civil war divides country north and south & Violence \\
               & 2010--2011 & Post-election crisis; Gbagbo refuses to concede & Violence, Regime \\
\midrule
Dem. Rep. & 1996--1997 & First Congo War; Mobutu overthrown & Violence, Regime \\
of the Congo & 1998--2003 & Second Congo War & Violence \\
             & 2018--12 & Contested election; Tshisekedi inaugurated & Within-Regime \\
\midrule
Eq. Guinea & 2004--03 & Failed coup attempt against Obiang & Regime \\
           & 2017--12 & Alleged mercenary plot; opposition crackdown & Violence, Within \\
\midrule
Gabon & 1993--12 & Disputed elections trigger riots & Protest, Within \\
      & 2016--08 & Post-election protests; parliament set ablaze & Protest, Violence \\
      & 2023--08 & Military coup removes President Ali Bongo & Regime \\
\midrule
Ghana & 2000--12 & Peaceful transfer from Rawlings to Kufuor & Within-Regime \\
      & 2008--12 & Peaceful alternation; opposition victory & Within-Regime \\
\midrule
Guinea-Bissau & 1998--1999 & Civil war; President Vieira ousted & Violence, Regime \\
              & 2003--09 & Military coup removes President Ialá & Regime \\
              & 2012--04 & Coup interrupts elections & Regime \\
\midrule
Mali & 1991--03 & Coup ends Traoré dictatorship & Regime \\
     & 2012--03 & Military coup amid Tuareg rebellion & Violence, Regime \\
     & 2020--08 & Coup ousts President Keïta & Regime \\
     & 2021--05 & Second coup; Goïta consolidates rule & Regime \\
\midrule
Mauritania & 2005--08 & Coup removes President Taya & Regime \\
            & 2008--08 & Coup ousts President Abdallahi & Regime \\Political Instability Article Counts by Country
            & 2019--08 & Peaceful transfer to Ghazouani & Within-Regime \\
\midrule
Niger & 1996--01 & Coup by Baré Maïnassara & Regime \\
       & 1999--04 & Maïnassara assassinated & Violence, Regime \\
       & 2010--02 & Coup removes President Tandja & Regime \\
       & 2023--07 & Military coup detains President Bazoum & Regime \\
\midrule
Nigeria & 1993--06 & Annulled elections trigger mass protests & Protest, Within \\
        & 2009--07 & Boko Haram uprising in Maiduguri & Violence \\
        & 2012--01 & Occupy Nigeria fuel subsidy protests & Protest \\
        & 2020--10 & \#EndSARS protests; Lekki shootings & Protest, Violence \\
        & 2024--07 & \#EndBadGovernance protests & Protest \\
\midrule
Senegal & 2011--06 & Y'en a Marre protests against Wade & Protest \\
         & 2021--2024 & Protests over Sall's third-term bid & Protest \\
         & 2025--02 & Casamance peace deal signed & Violence \\
\midrule
Togo & 1993--01 & Violent crackdown on pro-democracy protests & Violence, Protest \\
     & 2005--02 & Eyadéma dies; son installed; post-election violence & Violence, Regime \\
     & 2017--08 & Protests demand term limits & Protest \\
     & 2024--04 & Constitutional reform to parliamentary system & Within-Regime \\
\bottomrule
\end{tabular}
\label{tab:instability_events_detailed}
\vspace{0.5em}
\begin{minipage}{\textwidth}
{\footnotesize \flushleft \textit{Notes:} This table summarizes major political events in the sample countries and identifies which NBS dimension(s) are most likely related to the event. Dimension abbreviations: Violence = Politically Motivated Violence; Protest = Mass Civil Protest; Regime = Instability of the Regime; Within = Instability Within the Regime.}
\end{minipage}
\end{table}


% % LaTeX Table
% % Table for Political Violence
% % latex table generated in R 4.5.2 by xtable 1.8-4 package
% % Thu Dec 18 16:08:10 2025
% \begin{sidewaystable}[htbp]
% \scriptsize
% \centering
% \centering
% \caption{Correlation Matrix for Political Violence} 
% \label{tab:corr_Political Violence}
% \begin{tabular}{ll*{18}{r}}
%   \toprule
% Event Type & Source & BEN & BFA & CMR & CAF & TCD & CIV & COD & GNQ & GAB & GHA & GNB & MLI & MRT & NER & Nigeria & SEN & TGO & MEAN \\ 
%   \midrule
% Annual Coups & CSP & 0.01 & -0.01 & 0.01 & 0.23 & 0.14 & 0.02 & 0.04 & -0.02 & 0.24 & -0.00 & 0.03 & 0.02 & -0.00 & 0.12 & -0.07 &  & 0.01 & 0.05 \\ 
%   Anti-Gov. Demonst. & CNTS & -0.00 & 0.11 & 0.08 & 0.27 & -0.02 & 0.13 &  &  & 0.10 & 0.01 & -0.04 & 0.06 & -0.01 & -0.04 & 0.06 & 0.05 & 0.01 & 0.05 \\ 
%   Autocracy Index & CSP & -0.37 & -0.21 & 0.13 & -0.74 & -0.01 &  & -0.41 & 0.19 & -0.30 & 0.51 & 0.05 & -0.04 & 0.02 & -0.16 & 0.12 & 0.13 & -0.17 & -0.08 \\ 
%   Battles & ACLED & 0.25 & 0.48 & 0.25 & 0.61 & 0.52 & 0.73 & -0.03 & 0.26 & 0.41 & 0.16 & 0.54 & 0.42 & 0.36 & 0.32 & 0.04 & 0.19 & 0.15 & 0.33 \\ 
%   Civil Disorder & IRCG &  & 0.27 & 0.14 &  &  & 0.16 & -0.15 &  & -0.03 & -0.22 & -0.06 & 0.35 &  & 0.15 & 0.20 & -0.22 & -0.21 & 0.03 \\ 
%   Civil Unrest & BARRET &  & 0.71 &  &  & 0.24 &  & 0.48 &  &  & 0.15 &  & 0.58 & 0.24 & 0.33 & 0.60 & 0.40 & 0.32 & 0.40 \\ 
%   Civil War & IRCG &  & 0.29 & 0.20 &  &  & 0.22 & -0.11 &  & -0.04 & -0.22 & -0.06 & 0.41 &  & 0.11 & 0.33 & -0.23 & -0.26 & 0.05 \\ 
%   Control of Corrupt. & WWGI & 0.24 & 0.06 & 0.09 & -0.12 & 0.15 & -0.29 & -0.12 & -0.04 & 0.18 & 0.23 & 0.45 & -0.26 & -0.20 & 0.18 & 0.02 & 0.07 & -0.19 & 0.03 \\ 
%   Corruption Index & GALLUP & -0.37 & -0.44 & -0.50 & -0.08 & 0.30 & -0.11 & 0.15 &  & 0.10 & 0.11 &  & -0.04 & 0.44 & -0.18 & 0.32 & -0.13 & -0.04 & -0.03 \\ 
%   Coups & COUPS & -0.01 & 0.30 & -0.00 & 0.07 & 0.21 & 0.23 & -0.05 &  & 0.22 &  & 0.25 & 0.26 & 0.06 & 0.10 & -0.02 &  & 0.23 & 0.13 \\ 
%   Democracy Index & CSP & -0.37 & 0.42 & -0.13 & -0.73 & -0.02 &  & -0.40 &  & -0.26 & -0.49 & 0.09 & -0.04 & -0.12 & -0.15 & 0.24 & -0.15 & -0.00 & -0.14 \\ 
%   Election Index & GALLUP & 0.35 & 0.57 & 0.49 & 0.72 & -0.18 & -0.18 & -0.07 &  & -0.04 & -0.03 &  & 0.32 & -0.26 & 0.24 & 0.36 & 0.47 & 0.04 & 0.19 \\ 
%   Ethnic Viol. & UCDP &  &  & -0.01 & 0.27 & 0.04 & 0.03 & -0.02 &  &  &  &  & -0.02 &  & -0.02 & 0.03 & 0.06 & 0.01 & 0.04 \\ 
%   Explosions & ACLED & -0.12 & 0.34 & 0.31 & 0.39 & 0.19 & 0.34 & 0.03 &  &  &  & 0.28 & 0.28 & 0.05 & 0.02 & 0.16 & 0.19 & 0.49 & 0.21 \\ 
%   Gov. Crises & CNTS & 0.04 & 0.03 & -0.03 & 0.22 & -0.02 & -0.01 &  &  &  &  & 0.03 & 0.03 &  &  & 0.07 &  & 0.01 & 0.04 \\ 
%   Gov. Effect. & WWGI & 0.51 & -0.02 & 0.22 & -0.54 & -0.24 & -0.28 & -0.24 & -0.14 & -0.14 & -0.26 & 0.50 & -0.18 & -0.47 & 0.04 & -0.08 & -0.17 & -0.18 & -0.10 \\ 
%   Gov. Frac. & DPI & -0.70 & 0.12 & 0.18 & 0.22 & -0.16 & 0.30 & -0.31 & -0.17 & -0.08 & 0.43 & 0.49 & 0.34 & 0.02 & -0.24 & 0.21 & -0.03 & -0.27 & 0.02 \\ 
%   Gov. index & GALLUP & 0.30 & 0.52 & 0.45 & 0.43 & -0.04 & -0.25 & 0.03 &  & 0.02 & -0.02 &  & 0.10 & -0.16 & 0.33 & 0.08 & 0.46 & 0.15 & 0.16 \\ 
%   Instability of Regime &  & 0.59 & 0.54 & 0.37 & 0.66 & 0.81 & 0.58 & 0.69 & 0.37 & 0.46 & 0.42 & 0.57 & 0.62 & 0.41 & 0.43 & 0.22 & 0.50 & 0.69 & 0.53 \\ 
%   Instability within Regime &  & 0.73 & 0.65 & 0.69 & 0.74 & 0.79 & 0.78 & 0.62 & 0.52 & 0.81 & 0.51 & 0.74 & 0.79 & 0.64 & 0.63 & 0.42 & 0.73 & 0.89 & 0.69 \\ 
%   Internal Conflict & IRCG &  & 0.10 & 0.33 &  &  & 0.23 & 0.18 &  & -0.04 & 0.14 & -0.07 & 0.48 &  & 0.29 & 0.22 & 0.21 & 0.52 & 0.22 \\ 
%   Leadership Index & GALLUP & 0.31 & 0.52 & 0.27 & 0.17 & -0.30 & -0.16 & -0.14 &  & -0.06 & -0.03 &  & 0.22 & 0.16 & 0.30 & 0.25 & 0.43 & 0.40 & 0.16 \\ 
%   Legis. Frac. & DPI & -0.89 & 0.10 & -0.00 & 0.13 & -0.35 & 0.27 & -0.28 & -0.33 & -0.37 & -0.55 & 0.34 & 0.21 & -0.15 & -0.27 & 0.41 & -0.13 & -0.38 & -0.13 \\ 
%   Maj. Party & DPI & -0.19 & -0.00 & -0.10 & -0.23 & 0.20 & 0.05 & -0.30 & 0.21 & 0.30 & 0.37 & 0.10 & 0.32 & -0.65 & 0.14 & 0.34 & 0.07 & -0.59 & 0.00 \\ 
%   Mass Civil Protest &  & 0.52 & 0.60 & 0.58 & 0.41 & 0.65 & 0.79 & 0.59 & 0.53 & 0.72 & 0.37 & 0.71 & 0.60 & 0.54 & 0.47 & 0.59 & 0.67 & 0.79 & 0.60 \\ 
%   Nr. Seats Largest Party & CNTS & -0.03 & -0.02 & -0.01 & 0.09 & -0.02 & 0.04 &  & 0.08 & 0.08 & 0.13 & -0.03 & -0.02 & 0.05 & -0.00 & 0.12 & 0.09 & -0.01 & 0.03 \\ 
%   Party in Power & DPI & 0.65 & -0.08 & -0.48 & 0.12 & -0.36 & -0.32 & -0.02 & 0.44 & -0.08 & -0.32 & -0.35 & 0.17 & -0.38 & 0.31 & 0.61 & 0.21 & -0.29 & -0.01 \\ 
%   Pol. Relaxations & WHIV & -0.03 & -0.01 & 0.05 & 0.05 & 0.23 & 0.02 & 0.14 & -0.03 & 0.32 & 0.07 & -0.01 & 0.04 & 0.19 & 0.04 & -0.17 & 0.08 & 0.13 & 0.07 \\ 
%   Pol. Sanctions & WHIV & 0.02 & -0.05 & 0.16 & 0.06 & 0.02 & 0.03 & 0.13 & 0.46 & 0.35 & 0.09 & 0.14 & 0.12 & 0.23 & 0.08 & -0.10 & 0.02 & 0.30 & 0.12 \\ 
%   Pol. Stab. No Viol. & WWGI & -0.23 & -0.52 & -0.57 & -0.59 & 0.20 & -0.28 & 0.31 & 0.10 & -0.07 & 0.05 & 0.23 & -0.46 & -0.56 & -0.07 & -0.41 & 0.12 & 0.21 & -0.15 \\ 
%   Political Violence &  & 1.00 & 1.00 & 1.00 & 1.00 & 1.00 & 1.00 & 1.00 & 1.00 & 1.00 & 1.00 & 1.00 & 1.00 & 1.00 & 1.00 & 1.00 & 1.00 & 1.00 & 1.00 \\ 
%   Protest & WHIV & 0.19 & -0.07 & 0.11 & 0.02 & 0.10 & 0.20 & 0.26 & 0.23 & 0.23 & 0.17 & 0.09 & -0.03 & 0.20 & -0.02 & -0.02 & 0.06 & 0.40 & 0.12 \\ 
%   Protests & ACLED & 0.16 & 0.46 & 0.41 & 0.31 & 0.08 & 0.09 & -0.12 & 0.10 & 0.17 & 0.13 & 0.00 & 0.46 & -0.15 & 0.42 & 0.24 & 0.29 & 0.13 & 0.19 \\ 
%   Reg. Quality & WWGI & 0.09 & -0.42 & -0.01 & -0.36 & -0.21 & -0.25 & -0.24 & -0.22 & -0.24 & 0.19 & 0.11 & -0.21 & -0.16 & -0.23 & 0.33 & 0.39 & -0.03 & -0.09 \\ 
%   Riots & ACLED & 0.16 & 0.25 & 0.10 & 0.50 & 0.03 & 0.05 & -0.06 & -0.01 & 0.30 & 0.00 & -0.02 & 0.23 & 0.00 & 0.08 & 0.15 & 0.24 & 0.08 & 0.12 \\ 
%   Rule of Law & WWGI & -0.05 & 0.08 & 0.33 & -0.40 & 0.28 & -0.31 & -0.31 & -0.12 & -0.06 & 0.34 & 0.14 & -0.46 & -0.50 & 0.39 & 0.09 & 0.08 & -0.24 & -0.04 \\ 
%   Strat. Devel. & ACLED & 0.07 & 0.48 & -0.10 & 0.33 & 0.24 & -0.01 & -0.01 & 0.16 & 0.36 & 0.13 & 0.19 & 0.11 & -0.05 & -0.06 & -0.03 & 0.38 & 0.08 & 0.13 \\ 
%   Terror Attack & CSP &  &  &  &  &  &  &  &  &  &  &  &  &  &  &  &  &  &  \\ 
%   Terrorism & IRCG &  & 0.33 & 0.31 &  &  & 0.14 & -0.16 &  & -0.03 & -0.22 & -0.07 & 0.41 &  & 0.14 & 0.23 & -0.25 & -0.26 & 0.05 \\ 
%   Uncert. Index & WUI & -0.03 & 0.31 & 0.23 & 0.07 & 0.04 & 0.05 & -0.09 &  & 0.02 & -0.05 & -0.04 & 0.33 & -0.04 & 0.18 & 0.30 & -0.04 & 0.01 & 0.08 \\ 
%   Viol. & WHIV & 0.53 & 0.01 & 0.35 & 0.14 & 0.50 & 0.59 & 0.53 & 0.08 & 0.62 & 0.25 & 0.50 & 0.09 & 0.30 & 0.14 & -0.07 & 0.38 & 0.59 & 0.33 \\ 
%   Viol. against Civilians & ACLED & 0.08 & 0.51 & 0.19 & 0.63 & 0.22 & 0.59 & -0.09 & -0.02 & 0.28 & 0.10 & 0.25 & 0.29 & 0.11 & 0.15 & 0.05 & 0.13 & 0.05 & 0.21 \\ 
%   Voice and Account. & WWGI & -0.44 & 0.07 & -0.07 & -0.45 & -0.16 & -0.33 & -0.19 & -0.27 & -0.18 & 0.40 & 0.09 & -0.38 & -0.41 & -0.02 & 0.28 & 0.24 & -0.14 & -0.12 \\ 
%   Years in Office & DPI & 0.46 & 0.12 & -0.48 & 0.08 & -0.36 & -0.16 & -0.14 & 0.44 & 0.09 & 0.21 & -0.33 & 0.72 & -0.25 & -0.31 & -0.13 & 0.16 & -0.19 & -0.00 \\ 
%    \bottomrule
% \end{tabular}
% \end{sidewaystable}



% % LaTeX Table
% % Mean correlations across countries
% % latex table generated in R 4.5.2 by xtable 1.8-4 package
% % Thu Dec 18 16:18:46 2025
% \begin{table}[htbp]
% \scriptsize
% \centering
% \centering
% \caption{Mean Correlations across Countries} 
% \label{tab:mean_correlation_across_factors}
% \begin{tabular}{ll*{4}{r}}
%   \toprule
% Event Type & Source & Pol. Viol. & Prot. & Within & Regime \\ 
%   \midrule
% Annual Coups & CSP & 0.05 & 0.01 & 0.03 & 0.04 \\ 
%   Anti-Gov. Demonst. & CNTS & 0.05 & 0.01 & 0.02 & 0.04 \\ 
%   Autocracy Index & CSP & -0.08 & -0.11 & -0.11 & -0.13 \\ 
%   Battles & ACLED & 0.33 & 0.18 & 0.19 & 0.14 \\ 
%   Civil Disorder & IRCG & 0.03 & -0.06 & -0.04 & -0.08 \\ 
%   Civil Unrest & BARRET & 0.40 & 0.48 & 0.56 & 0.37 \\ 
%   Civil War & IRCG & 0.05 & -0.06 & -0.03 & -0.08 \\ 
%   Control of Corrupt. & WWGI & 0.03 & -0.06 & 0.00 & -0.03 \\ 
%   Corruption Index & GALLUP & -0.03 & 0.02 & -0.01 & -0.06 \\ 
%   Coups & COUPS & 0.13 & 0.13 & 0.15 & 0.15 \\ 
%   Democracy Index & CSP & -0.14 & -0.14 & -0.08 & -0.01 \\ 
%   Election Index & GALLUP & 0.19 & 0.12 & 0.17 & 0.20 \\ 
%   Ethnic Viol. & UCDP & 0.04 & 0.00 & 0.03 & 0.02 \\ 
%   Explosions & ACLED & 0.21 & 0.08 & 0.09 & 0.07 \\ 
%   Gov. Crises & CNTS & 0.04 & 0.00 & 0.02 & 0.06 \\ 
%   Gov. Effect. & WWGI & -0.10 & -0.09 & -0.10 & -0.12 \\ 
%   Gov. Frac. & DPI & 0.02 & -0.00 & -0.01 & -0.00 \\ 
%   Gov. index & GALLUP & 0.16 & 0.06 & 0.15 & 0.17 \\ 
%   Instability of Regime &  & 0.53 & 0.48 & 0.64 & 1.00 \\ 
%   Instability within Regime &  & 0.69 & 0.60 & 1.00 & 0.64 \\ 
%   Internal Conflict & IRCG & 0.22 & 0.18 & 0.17 & 0.15 \\ 
%   Leadership Index & GALLUP & 0.16 & 0.04 & 0.13 & 0.17 \\ 
%   Legis. Frac. & DPI & -0.13 & -0.20 & -0.17 & -0.15 \\ 
%   Maj. Party & DPI & 0.00 & 0.06 & -0.02 & -0.18 \\ 
%   Mass Civil Protest &  & 0.60 & 1.00 & 0.60 & 0.48 \\ 
%   Nr. Seats Largest Party & CNTS & 0.03 & 0.01 & 0.01 & 0.00 \\ 
%   Party in Power & DPI & -0.01 & -0.00 & -0.07 & -0.16 \\ 
%   Pol. Relaxations & WHIV & 0.07 & 0.13 & 0.08 & 0.07 \\ 
%   Pol. Sanctions & WHIV & 0.12 & 0.22 & 0.15 & 0.13 \\ 
%   Pol. Stab. No Viol. & WWGI & -0.15 & -0.05 & -0.01 & -0.02 \\ 
%   Political Violence &  & 1.00 & 0.60 & 0.69 & 0.53 \\ 
%   Protest & WHIV & 0.12 & 0.27 & 0.18 & 0.16 \\ 
%   Protests & ACLED & 0.19 & 0.18 & 0.16 & 0.19 \\ 
%   Reg. Quality & WWGI & -0.09 & -0.05 & -0.08 & -0.08 \\ 
%   Riots & ACLED & 0.12 & 0.11 & 0.09 & 0.08 \\ 
%   Rule of Law & WWGI & -0.04 & -0.08 & -0.04 & -0.08 \\ 
%   Strat. Devel. & ACLED & 0.13 & 0.06 & 0.15 & 0.14 \\ 
%   Terror Attack & CSP &  &  &  &  \\ 
%   Terrorism & IRCG & 0.05 & -0.07 & -0.05 & -0.08 \\ 
%   Uncert. Index & WUI & 0.08 & 0.02 & 0.05 & 0.03 \\ 
%   Viol. & WHIV & 0.33 & 0.33 & 0.26 & 0.22 \\ 
%   Viol. against Civilians & ACLED & 0.21 & 0.07 & 0.14 & 0.08 \\ 
%   Voice and Account. & WWGI & -0.12 & -0.11 & -0.10 & -0.12 \\ 
%   Years in Office & DPI & -0.00 & 0.03 & -0.03 & -0.12 \\ 
%    \bottomrule
% \end{tabular}
% \end{table}







\end{document}
